\chapter{Trình bày theo mẫu bài thi và đối chiếu ảnh}
\label{ch:presentation}
\section{Trật tự đọc một lời giải}
Mẫu ảnh người dùng gửi có một điểm đáng học: mỗi phép tính xuất hiện sau lý do phải tính nó.
Bài giải đi từ dữ liệu và mô hình, qua công thức tổng quát, bảng cần lập, hệ đã thế số, quá trình tính, rồi mới kết luận.
Trong phiên bản này, cùng dàn ý đó được dùng trong GUI, HTML xuất ra và danh mục 58 lựa chọn ở cuối phụ lục.
Các chương trước vẫn giữ công thức, bảng tổng quát, ví dụ và phần Gauss/QR/SVD đầy đủ.

\begin{enumerate}
\item \textbf{Dữ liệu:} số quan sát $N$, các cặp gốc, mục tiêu, số mốc, đơn vị, trọng số và điều kiện biên. Nếu chọn mốc, công bố ID nguồn và thứ tự sau chọn.
\item \textbf{Đặt bài:} mô hình/cơ sở, biến chuẩn hóa, nhánh nghiệm hay hệ số quy tắc. Tách tiền xử lý ra khỏi phép giải.
\item \textbf{Công thức trước, số sau:} định nghĩa mỗi ô bảng, phép truy hồi hoặc từng tổng; viết dạng tổng gọn và dạng khai triển để biết phải cộng gì.
\item \textbf{Bảng làm việc:} bảng sai phân/tổng/trọng số/hệ số/moment/tầng. Mỗi ô cần hiểu được nó lấy dữ liệu nào từ bước trước.
\item \textbf{Tính:} thế số vào hệ/công thức; với vòng lặp ghi trạng thái đầu, các lượng tính mới, quyết định cập nhật và điều kiện dừng.
\item \textbf{Trả lời:} viết hàm hay đa thức cuối, giá trị được hỏi, đơn vị, đổi về biến gốc, sai số và giới hạn của kết luận.
\end{enumerate}
Các bảng nhỏ cạnh nhau khi đọc được; bảng mốc nhỏ và Horner dùng chiều ngang.
Bảng rộng được chia cột khi in và lặp cột chỉ số, không tự cắt mất hệ số.
Trong GUI/HTML, nhật ký dài đặt trong Chi tiết, vẫn giữ đầy đủ và mở hết khi in HTML.
Màu nền hoặc ô đánh dấu dùng để chỉ ô thực sự tham gia phép tính, không thay lời giải thích.

\section{Những điểm không được sao chép từ mẫu}
\subsection{Chọn 102 điểm âm khác với giữ toàn bộ 104 điểm}
Với mô hình $y=Ae^{bx+c/x}$, $A$ có dấu cố định.
Lấy $Z_i=\ln(-y_i)$ chỉ hợp lệ khi $y_i<0$.
Bài mẫu loại hai quan sát dương; tool không tự thực hiện lựa chọn tập con này.
Nếu chủ động xử lý ở Excel, phải nêu các ID bị loại và tính RMSE trên đúng tập đã công bố.
Tệp \texttt{Tien\_xu\_ly\_Excel.pdf} hướng dẫn các thao tác đó, kèm các trường hợp tịnh tiến và đổi đơn vị.

Khớp $Z$ bằng bình phương tối thiểu không tối thiểu hóa SSE của $y$.
Không được gọi nhánh log đã lọc điểm và nhánh LM trên $y$ giữ toàn bộ điểm là cùng một bài giải số.
Đặt trọng số 0 cũng không bỏ qua kiểm tra miền log của những quan sát vẫn đang được nhập.

\subsection{Cách chọn gốc, bậc và điều kiện dừng nội suy ngược}
Mã giả trên ảnh tự chọn hướng, lấy gốc tại một đầu khoảng đổi dấu và giảm bậc theo số mốc phía đó.
Tool giữ hướng và tập mốc người dùng đã chọn; gốc Newton tiến/lùi là đầu/cuối của tập đó, khoảng nghiệm có thể nằm ở giữa.
Hai lựa chọn này có thể tạo hai đa thức và hai phép lặp khác nhau.

Điều kiện trên ảnh $|t_{k+1}-t_k|<\varepsilon$ là ngưỡng theo biến không thứ nguyên $t$.
Tool dùng ngưỡng theo $x$ cùng phần dư trên $y$; chế độ chứng nhận còn đòi miền bất biến và cận co.
Mẫu bỏ khoảng không hội tụ; tool giữ trạng thái chưa đạt và các vòng đã thực hiện.
Vì vậy chỉ áp dụng cách \emph{trình bày} vòng lặp, không tự thay tiêu chí dừng hay giấu nhánh thất bại.

Mã giả nói thường dùng bậc tối đa 3 hoặc 4, nhưng ví dụ sau đó dùng sáu mốc trên hai cửa sổ khác nhau.
Cần khai báo chính xác cửa sổ và bậc trước khi đối chiếu hai kết quả; không lấy hai nghiệm in trên ảnh làm chuẩn duy nhất khi thiếu bảng dữ liệu đầy đủ.

\subsection{Ký hiệu đếm và chữ số}
Dùng $N$ cho số dòng, $n=N-1$ cho chỉ số cuối nếu cần. Mẫu dùng $n$ theo hai nghĩa ở các đoạn khác nhau.
Giữ hệ số và số trung gian đầy đủ, chỉ làm tròn khi trình bày.
Hệ chuẩn có thể trình bày để suy luận; nếu nghiệm thực sự tính bằng QR/SVD thì nhật ký phải ghi QR/SVD, không dựng một quá trình khử Gauss giả.

\section{Viết bình phương tối thiểu theo trật tự trên ảnh}
Xét $N$ cặp hợp lệ $x_i\ne0,y_i<0$, $i=0,\ldots,N-1$, và họ
\[
y=Ae^{bx+c/x},\qquad A<0.
\]
\textbf{Bước 1 -- tuyến tính hóa và công bố ý nghĩa tham số.}
\begin{align*}
-y&=(-A)e^{bx+c/x},\\
\ln(-y)&=\ln(-A)+bx+\frac cx,\\
Z&=C_0\varphi_0(x)+b\varphi_1(x)+c\varphi_2(x),\qquad
(\varphi_0,\varphi_1,\varphi_2)=(1,x,1/x).
\end{align*}
Ở đây $C_0=\ln(-A)$; khôi phục $A=-e^{C_0}$.
Trong tool chọn Hàm mũ, nhập cơ sở \texttt{x;1/x} và cặp $x,y$ gốc hợp lệ.

\textbf{Bước 2 -- lập hàm mục tiêu.}
\begin{align*}
J(C_0,b,c)&=\sum_{i=0}^{N-1}\left(C_0+bx_i+\frac{c}{x_i}-Z_i\right)^2\\
&=\left(C_0+bx_0+c/x_0-Z_0\right)^2+\cdots\\
&\hspace{8mm}+\left(C_0+bx_{N-1}+c/x_{N-1}-Z_{N-1}\right)^2\longrightarrow\min.
\end{align*}
\textbf{Bước 3 -- lập bảng tổng quát.} Mỗi hàng tạo đủ các tích cần chuyển sang hệ chuẩn:
\begin{center}\small
\begin{tabular}{crrrrrrr}\toprule
Dòng&$x_i$&$x_i^2$&$1/x_i$&$1/x_i^2$&$Z_i$&$x_iZ_i$&$Z_i/x_i$\\\midrule
0&$x_0$&$x_0^2$&$1/x_0$&$1/x_0^2$&$Z_0$&$x_0Z_0$&$Z_0/x_0$\\
$i$&$x_i$&$x_i^2$&$1/x_i$&$1/x_i^2$&$Z_i$&$x_iZ_i$&$Z_i/x_i$\\
$N-1$&$x_{N-1}$&$x_{N-1}^2$&$1/x_{N-1}$&$1/x_{N-1}^2$&$Z_{N-1}$&$x_{N-1}Z_{N-1}$&$Z_{N-1}/x_{N-1}$\\\midrule
Tổng&$S_x$&$S_{xx}$&$S_{1/x}$&$S_{1/x^2}$&$S_Z$&$S_{xZ}$&$S_{Z/x}$\\\bottomrule
\end{tabular}\end{center}
Ví dụ $S_x=x_0+x_1+\cdots+x_{N-1}$ và
$S_{xZ}=x_0Z_0+x_1Z_1+\cdots+x_{N-1}Z_{N-1}$.

\textbf{Bước 4 -- từ cực trị đến hệ chuẩn.}
Lấy $\partial J/\partial C_0=\partial J/\partial b=\partial J/\partial c=0$:
\[
\begin{bmatrix}N&S_x&S_{1/x}\\S_x&S_{xx}&N\\S_{1/x}&N&S_{1/x^2}\end{bmatrix}
\begin{bmatrix}C_0\\b\\c\end{bmatrix}
=\begin{bmatrix}S_Z\\S_{xZ}\\S_{Z/x}\end{bmatrix}.
\]
Ô $N$ ở hàng 2 cột 3 xuất phát từ $\sum_i x_i(1/x_i)=\sum_i1=N$.
Nếu có trọng số thì mỗi tổng nhân $w_i$, các ô $N$ thay bằng $\sum w_i$.

\textbf{Bước 5 -- thế số công khai.} Chỉ dùng các tổng \emph{đã làm tròn trên ảnh}, ta có hệ
\[
\begin{bmatrix}
102&694.365000&22.591638\\
694.365000&5896.316975&102\\
22.591638&102&8.528590
\end{bmatrix}
\begin{bmatrix}C_0\\b\\c\end{bmatrix}
=\begin{bmatrix}-103.195072\\-1107.657050\\-3.786516\end{bmatrix}.
\]
Giải hệ này được xấp xỉ $C_0=0.8684645$, $b=-0.30595013$, $c=0.91461331$,
$A=-2.38324857$, khớp chữ số in của tham số trong mẫu.
Đây là \emph{kiểm tra hệ số từ các tổng đã in}, không phải chạy lại 102 quan sát gốc.

\textbf{Bước 6 -- trả lời và sai số.}
\[
g(x)=-2.38324857\exp\left(-0.30595013x+\frac{0.91461331}{x}\right).
\]
Muốn kiểm tra RMSE $0.111911$ trên ảnh cần đủ các quan sát của chính bảng test1.
Repo hiện có \texttt{20251GKtest2.xlsx}; không dùng bảng đó thay cho test1.
Khi có bảng đúng, tính $r_i=y_i-g(x_i)$ rồi
\[
\operatorname{RMSE}_y=\sqrt{(r_0^2+\cdots+r_{N-1}^2)/N}.
\]
Không nhầm RMSE này với $\sqrt{\mathrm{SSE}/(N-3)}$ hoặc mục tiêu log $J$.

\section{Kiểm tra độc lập đúng bảy mốc Stirling trên ảnh}
\textbf{Input.} Lấy đúng bảy cặp đọc được trên ảnh; ID trong bảng toàn bộ của ảnh là 24 đến 30, chỉ số cục bộ $j=-3,\ldots,3$.
Nếu nhập riêng bảy dòng vào tool, ID nguồn trong lượt chạy là 0 đến 6.
\begin{center}\small
\begin{tabular}{crrrrrrr}\toprule
$j$&$-3$&$-2$&$-1$&0&1&2&3\\\midrule
ID trên ảnh&24&25&26&27&28&29&30\\
$x_j$&3.760&3.875&3.990&4.105&4.220&4.335&4.450\\
$y_j$&$-.9384$&$-.8979$&$-.8110$&$-.8056$&$-.7548$&$-.7416$&$-.7049$\\\bottomrule
\end{tabular}\end{center}
\textbf{Bước 1 -- đổi biến.}
$h=0.115$, $x_0=4.105$, $t=(x-4.105)/0.115$.
Tại $\bar x=4.15$, $\bar t=9/23\simeq0.391304347826$.

\textbf{Bước 2 -- bảng sai phân.} Mỗi ô cấp sau là hiệu hai ô liền kề cấp trước; các ô trống không phải số 0:
\begin{center}\small
\begin{tabular}{crrrrrrr}\toprule
$j$&$y_j$&$\Delta$&$\Delta^2$&$\Delta^3$&$\Delta^4$&$\Delta^5$&$\Delta^6$\\\midrule
$-3$&$-.9384$&.0405&.0464&$-.1279$&.2548&$-.4647$&.8187\\
$-2$&$-.8979$&.0869&$-.0815$&.1269&$-.2099$&.3540&\\
$-1$&$-.8110$&.0054&.0454&$-.0830$&.1441&&\\
0&$-.8056$&.0508&$-.0376$&.0611&&&\\
1&$-.7548$&.0132&.0235&&&&\\
2&$-.7416$&.0367&&&&&\\
3&$-.7049$&&&&&&\\\bottomrule
\end{tabular}\end{center}
Ví dụ $\Delta^2y_{-1}=0.0508-0.0054=0.0454$;
$\Delta^3y_{-1}=-0.0376-0.0454=-0.0830$.

\textbf{Bước 3 -- lấy đủ các hạng Stirling.}
\begin{align*}
Q_6(t)={}&-.8056+\frac{.0054+.0508}{2}t+\frac{.0454}{2!}t^2\\
&+\frac{.1269-.0830}{2\cdot3!}t(t^2-1)
 +\frac{-.2099}{4!}t^2(t^2-1)\\
&+\frac{-.4647+.3540}{2\cdot5!}t(t^2-1)(t^2-4)
 +\frac{.8187}{6!}t^2(t^2-1)(t^2-4).
\end{align*}
Hệ số đi với cơ sở lần lượt là
\[
\left[-.8056,\ .0281,\ .0227,\ \frac{.0439}{12},\
\frac{-.2099}{24},\ \frac{-.1107}{240},\ \frac{.8187}{720}\right].
\]
Các phân số giữ đủ chữ số; không nhân/chia giai thừa lần thứ hai.

\textbf{Bước 4 -- khai triển tích rồi cộng theo cột.}
\begin{center}\small
\begin{tabular}{lrrrrrrr}\toprule
Cơ sở&$t^6$&$t^5$&$t^4$&$t^3$&$t^2$&$t$&1\\\midrule
1&0&0&0&0&0&0&1\\
$t$&0&0&0&0&0&1&0\\
$t^2$&0&0&0&0&1&0&0\\
$t(t^2-1)$&0&0&0&1&0&$-1$&0\\
$t^2(t^2-1)$&0&0&1&0&$-1$&0&0\\
$t(t^2-1)(t^2-4)$&0&1&0&$-5$&0&4&0\\
$t^2(t^2-1)(t^2-4)$&1&0&$-5$&0&4&0&0\\\bottomrule
\end{tabular}\end{center}
Nhân từng hàng với hệ số của chính cơ sở đó, cộng các cột. Ví dụ
$[t^4]Q_6=-.2099/24-5(.8187/720)=-.01443125$.
\begin{align*}
Q_6(t)={}&.00113708333333333t^6-.00046125t^5-.01443125t^4\\
&+.00596458333333333t^3+.0359941666666667t^2\\
&+.0225966666666667t-.8056.
\end{align*}
Các số thập phân tuần hoàn trên chỉ là bản hiển thị; phép tính dùng phân số đầy đủ.

\textbf{Bước 5 -- Horner, giá trị và đổi đạo hàm.}
Với hệ số bậc giảm $a_6,\ldots,a_0$ và $c=9/23$, khuôn Horner là
\begin{center}\small
\begin{tabular}{lrrrrr}\toprule
Hàng&$t^6$&$t^5$&$\cdots$&$t$&1\\\midrule
Hệ số&$a_6$&$a_5$&$\cdots$&$a_1$&$a_0$\\
Nhân $c$&&$cb_6$&$\cdots$&$cb_2$&$cb_1$\\
Cộng&$b_6=a_6$&$b_5=a_5+cb_6$&$\cdots$&$b_1$&$b_0=Q_6(c)$\\\bottomrule
\end{tabular}\end{center}
Bỏ số dư $b_0$, chia tiếp dãy thương $b_6,\ldots,b_1$ cho $t-c$;
số dư lượt hai là $Q_6'(c)$. Tool và HTML minh họa ghi đủ các ô số của hai lượt.
\[
Q_6(9/23)\simeq-0.791227553804,\qquad
Q_6'(9/23)\simeq0.050055749174.
\]
\[
f(4.15)\simeq-0.791227553804,\qquad
f'(4.15)\simeq\frac{Q_6'(9/23)}{0.115}\simeq0.435267384125.
\]
\textbf{Sai lệch của ảnh:} giá trị $-0.791128$ trong mẫu gần với tổng mới gồm ba hạng đầu
$-.8056+.0281t+.0227t^2\simeq-.7911285$, chưa phải giá trị của đa thức bậc 6 đã viết.
Đạo hàm trong ảnh khớp sau làm tròn. Không sửa kết quả đúng của tool để chạy theo giá trị thiếu hạng.

Tệp mở trực tiếp: \texttt{examples/presentation/Stirling\_7\_tu\_anh.json};
HTML cùng tên đọc ngoại tuyến. Đây là kiểm tra chính xác bảy mốc nhìn thấy, không xác nhận toàn bộ 104 mốc hay RMSE của dữ liệu chưa có.

\section{Đọc các bảng lặp đúng ý nghĩa}
\textbf{Nội suy ngược:} trạng thái $t_k$ $\to$ tích từng cơ sở $\to$ tổng phi tuyến $\to$ $H(t_k)$ $\to$ $x_{k+1}$ $\to$ độ đổi/phần dư/cận co $\to$ quyết định.
\textbf{LM:} tham số cũ $\to$ phần dư/Jacobian $\to$ damping/hệ tăng cường $\to$ bước thử $\to$ SSE mới $\to$ nhận hoặc từ chối.
\textbf{RK:} trạng thái cũ $\to$ thời gian/trạng thái từng tầng $\to$ độ dốc $K_j$ $\to$ đóng góp $hb_jK_j$ $\to$ nghiệm mới.
\textbf{Adams:} ghi riêng các bước RK4 khởi tạo, sau đó AB dự báo và AM hiệu chỉnh. Hệ số lịch sử được công bố đúng vị trí.
\textbf{Bước ODE ẩn:} công bố điểm đầu của bộ giải, nghiệm cuối, phần dư, số đánh giá hàm/Jacobian và thông báo thư viện. Số đánh giá không phải số vòng Newton.

Spline và nội suy đa thức là tính hữu hạn, không có tiêu chí hội tụ của phép lặp nghiệm.
Trị riêng và bài toán biên dùng bộ giải thư viện; không dựng nhật ký vòng lặp nội bộ khi thư viện không cung cấp.
Trị riêng trong tool chuẩn hóa $\max|u|=1$, không tự đổi sang chuẩn tích phân có trọng số.

\section{Dàn ý trình bày của đủ 58 lựa chọn}

Các dàn ý sau lấy từ cùng danh mục mà GUI và HTML sử dụng. Đây là đường đi của bài giải; các bảng tổng quát, phép khử, công thức và ví dụ chi tiết vẫn được giữ trong các chương trước. Không thay lời giải số bằng một đoạn mã giả.

Trong tool, nhóm Chi tiết có thể mở ra; bản HTML in ra mở toàn bộ chi tiết. Bản sao chép và các bảng xuất giữ mọi bước đã lưu.

\subsection{Nội suy đa thức}

\examplebreak\subsubsection*{1. Lagrange}

\textbf{Input:} Điểm cần tính x;\allowbreak{} Cấp đạo hàm (\allowbreak{}0 =\allowbreak{} giá trị hàm);\allowbreak{} Số mốc dùng (\allowbreak{}để trống: toàn bảng).

Bảng cặp $x,y$ là đầu vào riêng của lượt chạy; giữ ánh xạ nguồn nếu chọn một phần bảng.

\begin{enumerate}

\item Công bố tập mốc thực dùng và ID nguồn;\allowbreak{} kiểm tra x phân biệt,\allowbreak{} điều kiện lưới và số mốc của phương pháp.

\item Đặt biến và gốc: với mốc cách đều dùng t=\allowbreak{}(\allowbreak{}x\ensuremath{-}x\_gốc)/\allowbreak{}h;\allowbreak{} với Lagrange/\allowbreak{}Newton bất kỳ giữ biến x.

\item Lập các tích W\_i và mẫu d\_i;\allowbreak{} lấy hệ số y\_i/\allowbreak{}d\_i. Ghi rõ nguồn của từng hệ số và mẫu chia.

\item Viết đa thức dưới dạng tổng các hệ số nhân cơ sở;\allowbreak{} khai triển từng tích và cộng theo cột khi cần dạng chính tắc.

\item Thế điểm tính,\allowbreak{} cộng đủ các hạng;\allowbreak{} tính đạo hàm trên cùng đa thức. Nếu đổi biến,\allowbreak{} P\textasciicircum{}(\allowbreak{}r)(\allowbreak{}x)=\allowbreak{}Q\textasciicircum{}(\allowbreak{}r)(\allowbreak{}t)/\allowbreak{}h\textasciicircum{}r.

\item Thế lại các mốc và đối chiếu cách đánh giá;\allowbreak{} kết luận giá trị,\allowbreak{} đạo hàm và phạm vi sai số.

\end{enumerate}

\textbf{Output:} Tập mốc,\allowbreak{} bảng hệ số,\allowbreak{} đa thức,\allowbreak{} giá trị tại điểm tính và đạo hàm được yêu cầu;\allowbreak{} cận sai số chỉ khi có giả thiết bổ sung.

\examplebreak\subsubsection*{2. Newton}

\textbf{Input:} Loại mốc;\allowbreak{} Hướng nội suy;\allowbreak{} Điểm cần tính x;\allowbreak{} Cấp đạo hàm (\allowbreak{}0 =\allowbreak{} giá trị hàm);\allowbreak{} Số mốc dùng (\allowbreak{}để trống: toàn bảng).

Bảng cặp $x,y$ là đầu vào riêng của lượt chạy; giữ ánh xạ nguồn nếu chọn một phần bảng.

\begin{enumerate}

\item Công bố tập mốc thực dùng và ID nguồn;\allowbreak{} kiểm tra x phân biệt,\allowbreak{} điều kiện lưới và số mốc của phương pháp.

\item Đặt biến và gốc: với mốc cách đều dùng t=\allowbreak{}(\allowbreak{}x\ensuremath{-}x\_gốc)/\allowbreak{}h;\allowbreak{} với Lagrange/\allowbreak{}Newton bất kỳ giữ biến x.

\item Lập sai phân;\allowbreak{} lấy đường hệ số Newton tiến. Ghi rõ nguồn của từng hệ số và mẫu chia.

\item Viết đa thức dưới dạng tổng các hệ số nhân cơ sở;\allowbreak{} khai triển từng tích và cộng theo cột khi cần dạng chính tắc.

\item Thế điểm tính,\allowbreak{} cộng đủ các hạng;\allowbreak{} tính đạo hàm trên cùng đa thức. Nếu đổi biến,\allowbreak{} P\textasciicircum{}(\allowbreak{}r)(\allowbreak{}x)=\allowbreak{}Q\textasciicircum{}(\allowbreak{}r)(\allowbreak{}t)/\allowbreak{}h\textasciicircum{}r.

\item Thế lại các mốc và đối chiếu cách đánh giá;\allowbreak{} kết luận giá trị,\allowbreak{} đạo hàm và phạm vi sai số.

\end{enumerate}

\textbf{Output:} Tập mốc,\allowbreak{} bảng hệ số,\allowbreak{} đa thức,\allowbreak{} giá trị tại điểm tính và đạo hàm được yêu cầu;\allowbreak{} cận sai số chỉ khi có giả thiết bổ sung.

\examplebreak\subsubsection*{3. Gauss I}

\textbf{Input:} Điểm cần tính x;\allowbreak{} Cấp đạo hàm (\allowbreak{}0 =\allowbreak{} giá trị hàm);\allowbreak{} Số mốc dùng (\allowbreak{}để trống: toàn bảng).

Bảng cặp $x,y$ là đầu vào riêng của lượt chạy; giữ ánh xạ nguồn nếu chọn một phần bảng.

\begin{enumerate}

\item Công bố tập mốc thực dùng và ID nguồn;\allowbreak{} kiểm tra x phân biệt,\allowbreak{} điều kiện lưới và số mốc của phương pháp.

\item Đặt biến và gốc: với mốc cách đều dùng t=\allowbreak{}(\allowbreak{}x\ensuremath{-}x\_gốc)/\allowbreak{}h;\allowbreak{} với Lagrange/\allowbreak{}Newton bất kỳ giữ biến x.

\item Lập bảng sai phân;\allowbreak{} kết nạp mốc quanh tâm theo thứ tự 0,\allowbreak{} +\allowbreak{}1,\allowbreak{} \ensuremath{-}1,\allowbreak{} +\allowbreak{}2,\allowbreak{} \ensuremath{-}2,\allowbreak{} …. Ghi rõ nguồn của từng hệ số và mẫu chia.

\item Viết đa thức dưới dạng tổng các hệ số nhân cơ sở;\allowbreak{} khai triển từng tích và cộng theo cột khi cần dạng chính tắc.

\item Thế điểm tính,\allowbreak{} cộng đủ các hạng;\allowbreak{} tính đạo hàm trên cùng đa thức. Nếu đổi biến,\allowbreak{} P\textasciicircum{}(\allowbreak{}r)(\allowbreak{}x)=\allowbreak{}Q\textasciicircum{}(\allowbreak{}r)(\allowbreak{}t)/\allowbreak{}h\textasciicircum{}r.

\item Thế lại các mốc và đối chiếu cách đánh giá;\allowbreak{} kết luận giá trị,\allowbreak{} đạo hàm và phạm vi sai số.

\end{enumerate}

\textbf{Output:} Tập mốc,\allowbreak{} bảng hệ số,\allowbreak{} đa thức,\allowbreak{} giá trị tại điểm tính và đạo hàm được yêu cầu;\allowbreak{} cận sai số chỉ khi có giả thiết bổ sung.

\examplebreak\subsubsection*{4. Gauss II}

\textbf{Input:} Điểm cần tính x;\allowbreak{} Cấp đạo hàm (\allowbreak{}0 =\allowbreak{} giá trị hàm);\allowbreak{} Số mốc dùng (\allowbreak{}để trống: toàn bảng).

Bảng cặp $x,y$ là đầu vào riêng của lượt chạy; giữ ánh xạ nguồn nếu chọn một phần bảng.

\begin{enumerate}

\item Công bố tập mốc thực dùng và ID nguồn;\allowbreak{} kiểm tra x phân biệt,\allowbreak{} điều kiện lưới và số mốc của phương pháp.

\item Đặt biến và gốc: với mốc cách đều dùng t=\allowbreak{}(\allowbreak{}x\ensuremath{-}x\_gốc)/\allowbreak{}h;\allowbreak{} với Lagrange/\allowbreak{}Newton bất kỳ giữ biến x.

\item Lập bảng sai phân;\allowbreak{} kết nạp mốc quanh tâm theo thứ tự 0,\allowbreak{} \ensuremath{-}1,\allowbreak{} +\allowbreak{}1,\allowbreak{} \ensuremath{-}2,\allowbreak{} +\allowbreak{}2,\allowbreak{} …. Ghi rõ nguồn của từng hệ số và mẫu chia.

\item Viết đa thức dưới dạng tổng các hệ số nhân cơ sở;\allowbreak{} khai triển từng tích và cộng theo cột khi cần dạng chính tắc.

\item Thế điểm tính,\allowbreak{} cộng đủ các hạng;\allowbreak{} tính đạo hàm trên cùng đa thức. Nếu đổi biến,\allowbreak{} P\textasciicircum{}(\allowbreak{}r)(\allowbreak{}x)=\allowbreak{}Q\textasciicircum{}(\allowbreak{}r)(\allowbreak{}t)/\allowbreak{}h\textasciicircum{}r.

\item Thế lại các mốc và đối chiếu cách đánh giá;\allowbreak{} kết luận giá trị,\allowbreak{} đạo hàm và phạm vi sai số.

\end{enumerate}

\textbf{Output:} Tập mốc,\allowbreak{} bảng hệ số,\allowbreak{} đa thức,\allowbreak{} giá trị tại điểm tính và đạo hàm được yêu cầu;\allowbreak{} cận sai số chỉ khi có giả thiết bổ sung.

\examplebreak\subsubsection*{5. Stirling}

\textbf{Input:} Điểm cần tính x;\allowbreak{} Cấp đạo hàm (\allowbreak{}0 =\allowbreak{} giá trị hàm);\allowbreak{} Số mốc dùng (\allowbreak{}để trống: toàn bảng).

Bảng cặp $x,y$ là đầu vào riêng của lượt chạy; giữ ánh xạ nguồn nếu chọn một phần bảng.

\begin{enumerate}

\item Công bố tập mốc thực dùng và ID nguồn;\allowbreak{} kiểm tra x phân biệt,\allowbreak{} điều kiện lưới và số mốc của phương pháp.

\item Đặt biến và gốc: với mốc cách đều dùng t=\allowbreak{}(\allowbreak{}x\ensuremath{-}x\_gốc)/\allowbreak{}h;\allowbreak{} với Lagrange/\allowbreak{}Newton bất kỳ giữ biến x.

\item Lập bảng sai phân;\allowbreak{} cấp lẻ lấy trung bình hai ô kề tâm,\allowbreak{} cấp chẵn lấy một ô giữa. Ghi rõ nguồn của từng hệ số và mẫu chia.

\item Viết đa thức dưới dạng tổng các hệ số nhân cơ sở;\allowbreak{} khai triển từng tích và cộng theo cột khi cần dạng chính tắc.

\item Thế điểm tính,\allowbreak{} cộng đủ các hạng;\allowbreak{} tính đạo hàm trên cùng đa thức. Nếu đổi biến,\allowbreak{} P\textasciicircum{}(\allowbreak{}r)(\allowbreak{}x)=\allowbreak{}Q\textasciicircum{}(\allowbreak{}r)(\allowbreak{}t)/\allowbreak{}h\textasciicircum{}r.

\item Thế lại các mốc và đối chiếu cách đánh giá;\allowbreak{} kết luận giá trị,\allowbreak{} đạo hàm và phạm vi sai số.

\end{enumerate}

\textbf{Output:} Tập mốc,\allowbreak{} bảng hệ số,\allowbreak{} đa thức,\allowbreak{} giá trị tại điểm tính và đạo hàm được yêu cầu;\allowbreak{} cận sai số chỉ khi có giả thiết bổ sung.

\examplebreak\subsubsection*{6. Bessel}

\textbf{Input:} Điểm cần tính x;\allowbreak{} Cấp đạo hàm (\allowbreak{}0 =\allowbreak{} giá trị hàm);\allowbreak{} Số mốc dùng (\allowbreak{}để trống: toàn bảng).

Bảng cặp $x,y$ là đầu vào riêng của lượt chạy; giữ ánh xạ nguồn nếu chọn một phần bảng.

\begin{enumerate}

\item Công bố tập mốc thực dùng và ID nguồn;\allowbreak{} kiểm tra x phân biệt,\allowbreak{} điều kiện lưới và số mốc của phương pháp.

\item Đặt biến và gốc: với mốc cách đều dùng t=\allowbreak{}(\allowbreak{}x\ensuremath{-}x\_gốc)/\allowbreak{}h;\allowbreak{} với Lagrange/\allowbreak{}Newton bất kỳ giữ biến x.

\item Lập bảng sai phân;\allowbreak{} ghép hai mốc giữa,\allowbreak{} cấp chẵn lấy trung bình hai ô đối xứng. Ghi rõ nguồn của từng hệ số và mẫu chia.

\item Viết đa thức dưới dạng tổng các hệ số nhân cơ sở;\allowbreak{} khai triển từng tích và cộng theo cột khi cần dạng chính tắc.

\item Thế điểm tính,\allowbreak{} cộng đủ các hạng;\allowbreak{} tính đạo hàm trên cùng đa thức. Nếu đổi biến,\allowbreak{} P\textasciicircum{}(\allowbreak{}r)(\allowbreak{}x)=\allowbreak{}Q\textasciicircum{}(\allowbreak{}r)(\allowbreak{}t)/\allowbreak{}h\textasciicircum{}r.

\item Thế lại các mốc và đối chiếu cách đánh giá;\allowbreak{} kết luận giá trị,\allowbreak{} đạo hàm và phạm vi sai số.

\end{enumerate}

\textbf{Output:} Tập mốc,\allowbreak{} bảng hệ số,\allowbreak{} đa thức,\allowbreak{} giá trị tại điểm tính và đạo hàm được yêu cầu;\allowbreak{} cận sai số chỉ khi có giả thiết bổ sung.

\subsection{Spline}

\examplebreak\subsubsection*{7. Tuyến tính}

\textbf{Input:} Điểm cần tính x (\allowbreak{}có thể để trống);\allowbreak{} Cấp đạo hàm (\allowbreak{}0 =\allowbreak{} giá trị);\allowbreak{} Cận tích phân a (\allowbreak{}không bắt buộc);\allowbreak{} Cận tích phân b (\allowbreak{}không bắt buộc);\allowbreak{} Cách giải hệ biên;\allowbreak{} Mức trình bày.

Bảng cặp $x,y$ là đầu vào riêng của lượt chạy; giữ ánh xạ nguồn nếu chọn một phần bảng.

\begin{enumerate}

\item Sắp cặp theo x,\allowbreak{} giữ ID nguồn;\allowbreak{} đặt h\_i=\allowbreak{}x\_(\allowbreak{}i+\allowbreak{}1)\ensuremath{-}x\_i và \ensuremath{\delta}\_i=\allowbreak{}(\allowbreak{}y\_(\allowbreak{}i+\allowbreak{}1)\ensuremath{-}y\_i)/\allowbreak{}h\_i.

\item Công bố bậc spline và điều kiện biên: không cần điều kiện biên bổ sung.

\item Đặt u=\allowbreak{}x\ensuremath{-}x\_i trên từng đoạn;\allowbreak{} lập điều kiện nội suy,\allowbreak{} ghép đạo hàm đến cấp p\ensuremath{-}1 và điều kiện biên.

\item Thực hiện truy hồi hoặc giải đúng hệ đã lập. Với cubic moment,\allowbreak{} tìm M\_i=\allowbreak{}S\ensuremath{\prime\prime}(\allowbreak{}x\_i) rồi khôi phục từng mảnh.

\item Viết S\_i(\allowbreak{}x) cùng miền [x\_i,\allowbreak{}x\_(\allowbreak{}i+\allowbreak{}1)];\allowbreak{} xác định mảnh chứa điểm tính trước khi thế số,\allowbreak{} đạo hàm hoặc tích phân.

\item Kiểm tra giá trị tại mốc,\allowbreak{} độ khớp đạo hàm và điều kiện biên;\allowbreak{} kết luận đúng miền nội suy đã nhập.

\end{enumerate}

\textbf{Output:} Điều kiện biên,\allowbreak{} hệ/\allowbreak{}truy hồi,\allowbreak{} hệ số và công thức mọi mảnh,\allowbreak{} giá trị/\allowbreak{}đạo hàm/\allowbreak{}tích phân được yêu cầu,\allowbreak{} phần dư kiểm tra.

\examplebreak\subsubsection*{8. Bậc hai — biên trái}

\textbf{Input:} Điểm cần tính x (\allowbreak{}có thể để trống);\allowbreak{} Cấp đạo hàm (\allowbreak{}0 =\allowbreak{} giá trị);\allowbreak{} Giá trị đạo hàm tại biên trái;\allowbreak{} Cận tích phân a (\allowbreak{}không bắt buộc);\allowbreak{} Cận tích phân b (\allowbreak{}không bắt buộc);\allowbreak{} Cách giải hệ biên;\allowbreak{} Mức trình bày.

Bảng cặp $x,y$ là đầu vào riêng của lượt chạy; giữ ánh xạ nguồn nếu chọn một phần bảng.

\begin{enumerate}

\item Sắp cặp theo x,\allowbreak{} giữ ID nguồn;\allowbreak{} đặt h\_i=\allowbreak{}x\_(\allowbreak{}i+\allowbreak{}1)\ensuremath{-}x\_i và \ensuremath{\delta}\_i=\allowbreak{}(\allowbreak{}y\_(\allowbreak{}i+\allowbreak{}1)\ensuremath{-}y\_i)/\allowbreak{}h\_i.

\item Công bố bậc spline và điều kiện biên: một đạo hàm đầu trái;\allowbreak{} truyền đạo hàm sang phải.

\item Đặt u=\allowbreak{}x\ensuremath{-}x\_i trên từng đoạn;\allowbreak{} lập điều kiện nội suy,\allowbreak{} ghép đạo hàm đến cấp p\ensuremath{-}1 và điều kiện biên.

\item Thực hiện truy hồi hoặc giải đúng hệ đã lập. Với cubic moment,\allowbreak{} tìm M\_i=\allowbreak{}S\ensuremath{\prime\prime}(\allowbreak{}x\_i) rồi khôi phục từng mảnh.

\item Viết S\_i(\allowbreak{}x) cùng miền [x\_i,\allowbreak{}x\_(\allowbreak{}i+\allowbreak{}1)];\allowbreak{} xác định mảnh chứa điểm tính trước khi thế số,\allowbreak{} đạo hàm hoặc tích phân.

\item Kiểm tra giá trị tại mốc,\allowbreak{} độ khớp đạo hàm và điều kiện biên;\allowbreak{} kết luận đúng miền nội suy đã nhập.

\end{enumerate}

\textbf{Output:} Điều kiện biên,\allowbreak{} hệ/\allowbreak{}truy hồi,\allowbreak{} hệ số và công thức mọi mảnh,\allowbreak{} giá trị/\allowbreak{}đạo hàm/\allowbreak{}tích phân được yêu cầu,\allowbreak{} phần dư kiểm tra.

\examplebreak\subsubsection*{9. Bậc hai — biên phải}

\textbf{Input:} Điểm cần tính x (\allowbreak{}có thể để trống);\allowbreak{} Cấp đạo hàm (\allowbreak{}0 =\allowbreak{} giá trị);\allowbreak{} Giá trị đạo hàm tại biên phải;\allowbreak{} Cận tích phân a (\allowbreak{}không bắt buộc);\allowbreak{} Cận tích phân b (\allowbreak{}không bắt buộc);\allowbreak{} Cách giải hệ biên;\allowbreak{} Mức trình bày.

Bảng cặp $x,y$ là đầu vào riêng của lượt chạy; giữ ánh xạ nguồn nếu chọn một phần bảng.

\begin{enumerate}

\item Sắp cặp theo x,\allowbreak{} giữ ID nguồn;\allowbreak{} đặt h\_i=\allowbreak{}x\_(\allowbreak{}i+\allowbreak{}1)\ensuremath{-}x\_i và \ensuremath{\delta}\_i=\allowbreak{}(\allowbreak{}y\_(\allowbreak{}i+\allowbreak{}1)\ensuremath{-}y\_i)/\allowbreak{}h\_i.

\item Công bố bậc spline và điều kiện biên: một đạo hàm đầu phải;\allowbreak{} truyền đạo hàm sang trái.

\item Đặt u=\allowbreak{}x\ensuremath{-}x\_i trên từng đoạn;\allowbreak{} lập điều kiện nội suy,\allowbreak{} ghép đạo hàm đến cấp p\ensuremath{-}1 và điều kiện biên.

\item Thực hiện truy hồi hoặc giải đúng hệ đã lập. Với cubic moment,\allowbreak{} tìm M\_i=\allowbreak{}S\ensuremath{\prime\prime}(\allowbreak{}x\_i) rồi khôi phục từng mảnh.

\item Viết S\_i(\allowbreak{}x) cùng miền [x\_i,\allowbreak{}x\_(\allowbreak{}i+\allowbreak{}1)];\allowbreak{} xác định mảnh chứa điểm tính trước khi thế số,\allowbreak{} đạo hàm hoặc tích phân.

\item Kiểm tra giá trị tại mốc,\allowbreak{} độ khớp đạo hàm và điều kiện biên;\allowbreak{} kết luận đúng miền nội suy đã nhập.

\end{enumerate}

\textbf{Output:} Điều kiện biên,\allowbreak{} hệ/\allowbreak{}truy hồi,\allowbreak{} hệ số và công thức mọi mảnh,\allowbreak{} giá trị/\allowbreak{}đạo hàm/\allowbreak{}tích phân được yêu cầu,\allowbreak{} phần dư kiểm tra.

\examplebreak\subsubsection*{10. Cubic tự nhiên}

\textbf{Input:} Điểm cần tính x (\allowbreak{}có thể để trống);\allowbreak{} Cấp đạo hàm (\allowbreak{}0 =\allowbreak{} giá trị);\allowbreak{} Cận tích phân a (\allowbreak{}không bắt buộc);\allowbreak{} Cận tích phân b (\allowbreak{}không bắt buộc);\allowbreak{} Cách giải hệ biên;\allowbreak{} Mức trình bày.

Bảng cặp $x,y$ là đầu vào riêng của lượt chạy; giữ ánh xạ nguồn nếu chọn một phần bảng.

\begin{enumerate}

\item Sắp cặp theo x,\allowbreak{} giữ ID nguồn;\allowbreak{} đặt h\_i=\allowbreak{}x\_(\allowbreak{}i+\allowbreak{}1)\ensuremath{-}x\_i và \ensuremath{\delta}\_i=\allowbreak{}(\allowbreak{}y\_(\allowbreak{}i+\allowbreak{}1)\ensuremath{-}y\_i)/\allowbreak{}h\_i.

\item Công bố bậc spline và điều kiện biên: S\ensuremath{\prime\prime}(\allowbreak{}a)=\allowbreak{}S\ensuremath{\prime\prime}(\allowbreak{}b)=\allowbreak{}0.

\item Đặt u=\allowbreak{}x\ensuremath{-}x\_i trên từng đoạn;\allowbreak{} lập điều kiện nội suy,\allowbreak{} ghép đạo hàm đến cấp p\ensuremath{-}1 và điều kiện biên.

\item Thực hiện truy hồi hoặc giải đúng hệ đã lập. Với cubic moment,\allowbreak{} tìm M\_i=\allowbreak{}S\ensuremath{\prime\prime}(\allowbreak{}x\_i) rồi khôi phục từng mảnh.

\item Viết S\_i(\allowbreak{}x) cùng miền [x\_i,\allowbreak{}x\_(\allowbreak{}i+\allowbreak{}1)];\allowbreak{} xác định mảnh chứa điểm tính trước khi thế số,\allowbreak{} đạo hàm hoặc tích phân.

\item Kiểm tra giá trị tại mốc,\allowbreak{} độ khớp đạo hàm và điều kiện biên;\allowbreak{} kết luận đúng miền nội suy đã nhập.

\end{enumerate}

\textbf{Output:} Điều kiện biên,\allowbreak{} hệ/\allowbreak{}truy hồi,\allowbreak{} hệ số và công thức mọi mảnh,\allowbreak{} giá trị/\allowbreak{}đạo hàm/\allowbreak{}tích phân được yêu cầu,\allowbreak{} phần dư kiểm tra.

\examplebreak\subsubsection*{11. Cubic kẹp}

\textbf{Input:} Điểm cần tính x (\allowbreak{}có thể để trống);\allowbreak{} Cấp đạo hàm (\allowbreak{}0 =\allowbreak{} giá trị);\allowbreak{} Giá trị đạo hàm tại biên trái;\allowbreak{} Giá trị đạo hàm tại biên phải;\allowbreak{} Cận tích phân a (\allowbreak{}không bắt buộc);\allowbreak{} Cận tích phân b (\allowbreak{}không bắt buộc);\allowbreak{} Cách giải hệ biên;\allowbreak{} Mức trình bày.

Bảng cặp $x,y$ là đầu vào riêng của lượt chạy; giữ ánh xạ nguồn nếu chọn một phần bảng.

\begin{enumerate}

\item Sắp cặp theo x,\allowbreak{} giữ ID nguồn;\allowbreak{} đặt h\_i=\allowbreak{}x\_(\allowbreak{}i+\allowbreak{}1)\ensuremath{-}x\_i và \ensuremath{\delta}\_i=\allowbreak{}(\allowbreak{}y\_(\allowbreak{}i+\allowbreak{}1)\ensuremath{-}y\_i)/\allowbreak{}h\_i.

\item Công bố bậc spline và điều kiện biên: hai giá trị S\ensuremath{\prime}(\allowbreak{}a),\allowbreak{} S\ensuremath{\prime}(\allowbreak{}b) được nhập.

\item Đặt u=\allowbreak{}x\ensuremath{-}x\_i trên từng đoạn;\allowbreak{} lập điều kiện nội suy,\allowbreak{} ghép đạo hàm đến cấp p\ensuremath{-}1 và điều kiện biên.

\item Thực hiện truy hồi hoặc giải đúng hệ đã lập. Với cubic moment,\allowbreak{} tìm M\_i=\allowbreak{}S\ensuremath{\prime\prime}(\allowbreak{}x\_i) rồi khôi phục từng mảnh.

\item Viết S\_i(\allowbreak{}x) cùng miền [x\_i,\allowbreak{}x\_(\allowbreak{}i+\allowbreak{}1)];\allowbreak{} xác định mảnh chứa điểm tính trước khi thế số,\allowbreak{} đạo hàm hoặc tích phân.

\item Kiểm tra giá trị tại mốc,\allowbreak{} độ khớp đạo hàm và điều kiện biên;\allowbreak{} kết luận đúng miền nội suy đã nhập.

\end{enumerate}

\textbf{Output:} Điều kiện biên,\allowbreak{} hệ/\allowbreak{}truy hồi,\allowbreak{} hệ số và công thức mọi mảnh,\allowbreak{} giá trị/\allowbreak{}đạo hàm/\allowbreak{}tích phân được yêu cầu,\allowbreak{} phần dư kiểm tra.

\examplebreak\subsubsection*{12. Cubic đạo hàm bậc hai}

\textbf{Input:} Điểm cần tính x (\allowbreak{}có thể để trống);\allowbreak{} Cấp đạo hàm (\allowbreak{}0 =\allowbreak{} giá trị);\allowbreak{} Giá trị đạo hàm tại biên trái;\allowbreak{} Giá trị đạo hàm tại biên phải;\allowbreak{} Cận tích phân a (\allowbreak{}không bắt buộc);\allowbreak{} Cận tích phân b (\allowbreak{}không bắt buộc);\allowbreak{} Cách giải hệ biên;\allowbreak{} Mức trình bày.

Bảng cặp $x,y$ là đầu vào riêng của lượt chạy; giữ ánh xạ nguồn nếu chọn một phần bảng.

\begin{enumerate}

\item Sắp cặp theo x,\allowbreak{} giữ ID nguồn;\allowbreak{} đặt h\_i=\allowbreak{}x\_(\allowbreak{}i+\allowbreak{}1)\ensuremath{-}x\_i và \ensuremath{\delta}\_i=\allowbreak{}(\allowbreak{}y\_(\allowbreak{}i+\allowbreak{}1)\ensuremath{-}y\_i)/\allowbreak{}h\_i.

\item Công bố bậc spline và điều kiện biên: hai giá trị S\ensuremath{\prime\prime}(\allowbreak{}a),\allowbreak{} S\ensuremath{\prime\prime}(\allowbreak{}b) được nhập.

\item Đặt u=\allowbreak{}x\ensuremath{-}x\_i trên từng đoạn;\allowbreak{} lập điều kiện nội suy,\allowbreak{} ghép đạo hàm đến cấp p\ensuremath{-}1 và điều kiện biên.

\item Thực hiện truy hồi hoặc giải đúng hệ đã lập. Với cubic moment,\allowbreak{} tìm M\_i=\allowbreak{}S\ensuremath{\prime\prime}(\allowbreak{}x\_i) rồi khôi phục từng mảnh.

\item Viết S\_i(\allowbreak{}x) cùng miền [x\_i,\allowbreak{}x\_(\allowbreak{}i+\allowbreak{}1)];\allowbreak{} xác định mảnh chứa điểm tính trước khi thế số,\allowbreak{} đạo hàm hoặc tích phân.

\item Kiểm tra giá trị tại mốc,\allowbreak{} độ khớp đạo hàm và điều kiện biên;\allowbreak{} kết luận đúng miền nội suy đã nhập.

\end{enumerate}

\textbf{Output:} Điều kiện biên,\allowbreak{} hệ/\allowbreak{}truy hồi,\allowbreak{} hệ số và công thức mọi mảnh,\allowbreak{} giá trị/\allowbreak{}đạo hàm/\allowbreak{}tích phân được yêu cầu,\allowbreak{} phần dư kiểm tra.

\examplebreak\subsubsection*{13. Cubic hỗn hợp}

\textbf{Input:} Điểm cần tính x (\allowbreak{}có thể để trống);\allowbreak{} Cấp đạo hàm (\allowbreak{}0 =\allowbreak{} giá trị);\allowbreak{} Cấp đạo hàm biên trái;\allowbreak{} Cấp đạo hàm biên phải;\allowbreak{} Giá trị tại biên trái;\allowbreak{} Giá trị tại biên phải;\allowbreak{} Cận tích phân a (\allowbreak{}không bắt buộc);\allowbreak{} Cận tích phân b (\allowbreak{}không bắt buộc);\allowbreak{} Cách giải hệ biên;\allowbreak{} Mức trình bày.

Bảng cặp $x,y$ là đầu vào riêng của lượt chạy; giữ ánh xạ nguồn nếu chọn một phần bảng.

\begin{enumerate}

\item Sắp cặp theo x,\allowbreak{} giữ ID nguồn;\allowbreak{} đặt h\_i=\allowbreak{}x\_(\allowbreak{}i+\allowbreak{}1)\ensuremath{-}x\_i và \ensuremath{\delta}\_i=\allowbreak{}(\allowbreak{}y\_(\allowbreak{}i+\allowbreak{}1)\ensuremath{-}y\_i)/\allowbreak{}h\_i.

\item Công bố bậc spline và điều kiện biên: đúng phía,\allowbreak{} cấp và giá trị đạo hàm đã chọn ở mỗi đầu.

\item Đặt u=\allowbreak{}x\ensuremath{-}x\_i trên từng đoạn;\allowbreak{} lập điều kiện nội suy,\allowbreak{} ghép đạo hàm đến cấp p\ensuremath{-}1 và điều kiện biên.

\item Thực hiện truy hồi hoặc giải đúng hệ đã lập. Với cubic moment,\allowbreak{} tìm M\_i=\allowbreak{}S\ensuremath{\prime\prime}(\allowbreak{}x\_i) rồi khôi phục từng mảnh.

\item Viết S\_i(\allowbreak{}x) cùng miền [x\_i,\allowbreak{}x\_(\allowbreak{}i+\allowbreak{}1)];\allowbreak{} xác định mảnh chứa điểm tính trước khi thế số,\allowbreak{} đạo hàm hoặc tích phân.

\item Kiểm tra giá trị tại mốc,\allowbreak{} độ khớp đạo hàm và điều kiện biên;\allowbreak{} kết luận đúng miền nội suy đã nhập.

\end{enumerate}

\textbf{Output:} Điều kiện biên,\allowbreak{} hệ/\allowbreak{}truy hồi,\allowbreak{} hệ số và công thức mọi mảnh,\allowbreak{} giá trị/\allowbreak{}đạo hàm/\allowbreak{}tích phân được yêu cầu,\allowbreak{} phần dư kiểm tra.

\examplebreak\subsubsection*{14. Cubic tuần hoàn}

\textbf{Input:} Điểm cần tính x (\allowbreak{}có thể để trống);\allowbreak{} Cấp đạo hàm (\allowbreak{}0 =\allowbreak{} giá trị);\allowbreak{} Cận tích phân a (\allowbreak{}không bắt buộc);\allowbreak{} Cận tích phân b (\allowbreak{}không bắt buộc);\allowbreak{} Cách giải hệ biên;\allowbreak{} Mức trình bày.

Bảng cặp $x,y$ là đầu vào riêng của lượt chạy; giữ ánh xạ nguồn nếu chọn một phần bảng.

\begin{enumerate}

\item Sắp cặp theo x,\allowbreak{} giữ ID nguồn;\allowbreak{} đặt h\_i=\allowbreak{}x\_(\allowbreak{}i+\allowbreak{}1)\ensuremath{-}x\_i và \ensuremath{\delta}\_i=\allowbreak{}(\allowbreak{}y\_(\allowbreak{}i+\allowbreak{}1)\ensuremath{-}y\_i)/\allowbreak{}h\_i.

\item Công bố bậc spline và điều kiện biên: giá trị hai đầu bằng nhau;\allowbreak{} ghép S\ensuremath{\prime} và S\ensuremath{\prime\prime} ở hai đầu.

\item Đặt u=\allowbreak{}x\ensuremath{-}x\_i trên từng đoạn;\allowbreak{} lập điều kiện nội suy,\allowbreak{} ghép đạo hàm đến cấp p\ensuremath{-}1 và điều kiện biên.

\item Thực hiện truy hồi hoặc giải đúng hệ đã lập. Với cubic moment,\allowbreak{} tìm M\_i=\allowbreak{}S\ensuremath{\prime\prime}(\allowbreak{}x\_i) rồi khôi phục từng mảnh.

\item Viết S\_i(\allowbreak{}x) cùng miền [x\_i,\allowbreak{}x\_(\allowbreak{}i+\allowbreak{}1)];\allowbreak{} xác định mảnh chứa điểm tính trước khi thế số,\allowbreak{} đạo hàm hoặc tích phân.

\item Kiểm tra giá trị tại mốc,\allowbreak{} độ khớp đạo hàm và điều kiện biên;\allowbreak{} kết luận đúng miền nội suy đã nhập.

\end{enumerate}

\textbf{Output:} Điều kiện biên,\allowbreak{} hệ/\allowbreak{}truy hồi,\allowbreak{} hệ số và công thức mọi mảnh,\allowbreak{} giá trị/\allowbreak{}đạo hàm/\allowbreak{}tích phân được yêu cầu,\allowbreak{} phần dư kiểm tra.

\examplebreak\subsubsection*{15. Cubic not-a-knot}

\textbf{Input:} Điểm cần tính x (\allowbreak{}có thể để trống);\allowbreak{} Cấp đạo hàm (\allowbreak{}0 =\allowbreak{} giá trị);\allowbreak{} Cận tích phân a (\allowbreak{}không bắt buộc);\allowbreak{} Cận tích phân b (\allowbreak{}không bắt buộc);\allowbreak{} Cách giải hệ biên;\allowbreak{} Mức trình bày.

Bảng cặp $x,y$ là đầu vào riêng của lượt chạy; giữ ánh xạ nguồn nếu chọn một phần bảng.

\begin{enumerate}

\item Sắp cặp theo x,\allowbreak{} giữ ID nguồn;\allowbreak{} đặt h\_i=\allowbreak{}x\_(\allowbreak{}i+\allowbreak{}1)\ensuremath{-}x\_i và \ensuremath{\delta}\_i=\allowbreak{}(\allowbreak{}y\_(\allowbreak{}i+\allowbreak{}1)\ensuremath{-}y\_i)/\allowbreak{}h\_i.

\item Công bố bậc spline và điều kiện biên: đạo hàm bậc ba liên tục tại mốc thứ hai và áp chót;\allowbreak{} xử lý riêng bảng ít mốc.

\item Đặt u=\allowbreak{}x\ensuremath{-}x\_i trên từng đoạn;\allowbreak{} lập điều kiện nội suy,\allowbreak{} ghép đạo hàm đến cấp p\ensuremath{-}1 và điều kiện biên.

\item Thực hiện truy hồi hoặc giải đúng hệ đã lập. Với cubic moment,\allowbreak{} tìm M\_i=\allowbreak{}S\ensuremath{\prime\prime}(\allowbreak{}x\_i) rồi khôi phục từng mảnh.

\item Viết S\_i(\allowbreak{}x) cùng miền [x\_i,\allowbreak{}x\_(\allowbreak{}i+\allowbreak{}1)];\allowbreak{} xác định mảnh chứa điểm tính trước khi thế số,\allowbreak{} đạo hàm hoặc tích phân.

\item Kiểm tra giá trị tại mốc,\allowbreak{} độ khớp đạo hàm và điều kiện biên;\allowbreak{} kết luận đúng miền nội suy đã nhập.

\end{enumerate}

\textbf{Output:} Điều kiện biên,\allowbreak{} hệ/\allowbreak{}truy hồi,\allowbreak{} hệ số và công thức mọi mảnh,\allowbreak{} giá trị/\allowbreak{}đạo hàm/\allowbreak{}tích phân được yêu cầu,\allowbreak{} phần dư kiểm tra.

\examplebreak\subsubsection*{16. Bậc 4 — đạo hàm trái}

\textbf{Input:} Điểm cần tính x (\allowbreak{}có thể để trống);\allowbreak{} Cấp đạo hàm (\allowbreak{}0 =\allowbreak{} giá trị);\allowbreak{} S\ensuremath{\prime} S\ensuremath{\prime}\ensuremath{\prime} S\ensuremath{\prime}\ensuremath{\prime}\ensuremath{\prime} ở đầu trái;\allowbreak{} Cận tích phân a (\allowbreak{}không bắt buộc);\allowbreak{} Cận tích phân b (\allowbreak{}không bắt buộc);\allowbreak{} Cách giải hệ biên;\allowbreak{} Mức trình bày.

Bảng cặp $x,y$ là đầu vào riêng của lượt chạy; giữ ánh xạ nguồn nếu chọn một phần bảng.

\begin{enumerate}

\item Sắp cặp theo x,\allowbreak{} giữ ID nguồn;\allowbreak{} đặt h\_i=\allowbreak{}x\_(\allowbreak{}i+\allowbreak{}1)\ensuremath{-}x\_i và \ensuremath{\delta}\_i=\allowbreak{}(\allowbreak{}y\_(\allowbreak{}i+\allowbreak{}1)\ensuremath{-}y\_i)/\allowbreak{}h\_i.

\item Công bố bậc spline và điều kiện biên: ba giá trị S\ensuremath{\prime},\allowbreak{} S\ensuremath{\prime\prime},\allowbreak{} S\ensuremath{\prime\prime\prime} ở đầu trái;\allowbreak{} truyền các đạo hàm sang phải.

\item Đặt u=\allowbreak{}x\ensuremath{-}x\_i trên từng đoạn;\allowbreak{} lập điều kiện nội suy,\allowbreak{} ghép đạo hàm đến cấp p\ensuremath{-}1 và điều kiện biên.

\item Thực hiện truy hồi hoặc giải đúng hệ đã lập. Với cubic moment,\allowbreak{} tìm M\_i=\allowbreak{}S\ensuremath{\prime\prime}(\allowbreak{}x\_i) rồi khôi phục từng mảnh.

\item Viết S\_i(\allowbreak{}x) cùng miền [x\_i,\allowbreak{}x\_(\allowbreak{}i+\allowbreak{}1)];\allowbreak{} xác định mảnh chứa điểm tính trước khi thế số,\allowbreak{} đạo hàm hoặc tích phân.

\item Kiểm tra giá trị tại mốc,\allowbreak{} độ khớp đạo hàm và điều kiện biên;\allowbreak{} kết luận đúng miền nội suy đã nhập.

\end{enumerate}

\textbf{Output:} Điều kiện biên,\allowbreak{} hệ/\allowbreak{}truy hồi,\allowbreak{} hệ số và công thức mọi mảnh,\allowbreak{} giá trị/\allowbreak{}đạo hàm/\allowbreak{}tích phân được yêu cầu,\allowbreak{} phần dư kiểm tra.

\examplebreak\subsubsection*{17. Biên tổng quát}

\textbf{Input:} Điểm cần tính x (\allowbreak{}có thể để trống);\allowbreak{} Cấp đạo hàm (\allowbreak{}0 =\allowbreak{} giá trị);\allowbreak{} Bậc spline;\allowbreak{} Điều kiện biên (\allowbreak{}JSON);\allowbreak{} Cận tích phân a (\allowbreak{}không bắt buộc);\allowbreak{} Cận tích phân b (\allowbreak{}không bắt buộc);\allowbreak{} Cách giải hệ biên;\allowbreak{} Mức trình bày.

Bảng cặp $x,y$ là đầu vào riêng của lượt chạy; giữ ánh xạ nguồn nếu chọn một phần bảng.

\begin{enumerate}

\item Sắp cặp theo x,\allowbreak{} giữ ID nguồn;\allowbreak{} đặt h\_i=\allowbreak{}x\_(\allowbreak{}i+\allowbreak{}1)\ensuremath{-}x\_i và \ensuremath{\delta}\_i=\allowbreak{}(\allowbreak{}y\_(\allowbreak{}i+\allowbreak{}1)\ensuremath{-}y\_i)/\allowbreak{}h\_i.

\item Công bố bậc spline và điều kiện biên: p\ensuremath{-}1 điều kiện độc lập đúng phía và cấp trong JSON;\allowbreak{} không tự thêm biên tự nhiên.

\item Đặt u=\allowbreak{}x\ensuremath{-}x\_i trên từng đoạn;\allowbreak{} lập điều kiện nội suy,\allowbreak{} ghép đạo hàm đến cấp p\ensuremath{-}1 và điều kiện biên.

\item Thực hiện truy hồi hoặc giải đúng hệ đã lập. Với cubic moment,\allowbreak{} tìm M\_i=\allowbreak{}S\ensuremath{\prime\prime}(\allowbreak{}x\_i) rồi khôi phục từng mảnh.

\item Viết S\_i(\allowbreak{}x) cùng miền [x\_i,\allowbreak{}x\_(\allowbreak{}i+\allowbreak{}1)];\allowbreak{} xác định mảnh chứa điểm tính trước khi thế số,\allowbreak{} đạo hàm hoặc tích phân.

\item Kiểm tra giá trị tại mốc,\allowbreak{} độ khớp đạo hàm và điều kiện biên;\allowbreak{} kết luận đúng miền nội suy đã nhập.

\end{enumerate}

\textbf{Output:} Điều kiện biên,\allowbreak{} hệ/\allowbreak{}truy hồi,\allowbreak{} hệ số và công thức mọi mảnh,\allowbreak{} giá trị/\allowbreak{}đạo hàm/\allowbreak{}tích phân được yêu cầu,\allowbreak{} phần dư kiểm tra.

\subsection{Bình phương tối thiểu}

\examplebreak\subsubsection*{18. Tuyến tính theo cơ sở}

\textbf{Input:} Các hàm cơ sở (\allowbreak{}ngăn bằng dấu ;\allowbreak{});\allowbreak{} Trọng số w\_i (\allowbreak{}để trống: tất cả bằng 1);\allowbreak{} Cách giải bình phương tối thiểu;\allowbreak{} Quy tắc khi thiếu hạng;\allowbreak{} Điểm dự đoán x (\allowbreak{}không bắt buộc);\allowbreak{} Mức trình bày.

Bảng cặp $x,y$ là đầu vào riêng của lượt chạy; giữ ánh xạ nguồn nếu chọn một phần bảng.

\begin{enumerate}

\item Công bố N quan sát,\allowbreak{} trọng số và mô hình. Kiểm tra miền,\allowbreak{} dấu và hạng;\allowbreak{} không tự lọc dòng để ép phép biến đổi hợp lệ.

\item Viết phép biến đổi và ánh xạ tham số: z=\allowbreak{}y,\allowbreak{} g(\allowbreak{}x)=\allowbreak{}c\_0\ensuremath{\varphi}\_0(\allowbreak{}x)+\allowbreak{}…+\allowbreak{}c\_(\allowbreak{}p\ensuremath{-}1)\ensuremath{\varphi}\_(\allowbreak{}p\ensuremath{-}1)(\allowbreak{}x);\allowbreak{} các cơ sở lấy đúng ô nhập.

\item Lập \ensuremath{\Phi}\_ij=\allowbreak{}\ensuremath{\varphi}\_j(\allowbreak{}x\_i),\allowbreak{} bảng các tích và các tổng;\allowbreak{} viết J=\allowbreak{}\ensuremath{\sum}w\_i(\allowbreak{}\ensuremath{\sum}c\_j\ensuremath{\varphi}\_j(\allowbreak{}x\_i)\ensuremath{-}z\_i)².

\item Từ \ensuremath{\partial}J/\allowbreak{}\ensuremath{\partial}c\_j=\allowbreak{}0 suy ra (\allowbreak{}\ensuremath{\Phi}\ensuremath{{}^T}W\ensuremath{\Phi})c=\allowbreak{}\ensuremath{\Phi}\ensuremath{{}^T}Wz. Trình bày hệ tổng quát rồi hệ đã thế số.

\item Tính nghiệm bằng bộ giải đã chọn: Gauss trên hệ chuẩn,\allowbreak{} QR trên thiết kế có trọng số hoặc SVD. Nhật ký phải đúng bộ giải.

\item Khôi phục tham số và hàm g theo biến gốc;\allowbreak{} lập r\_i=\allowbreak{}y\_i\ensuremath{-}g(\allowbreak{}x\_i),\allowbreak{} tính SSE,\allowbreak{} RMSE và mục tiêu trên thang khớp riêng biệt.

\end{enumerate}

\textbf{Output:} Hàm g(\allowbreak{}x),\allowbreak{} tham số gốc,\allowbreak{} bảng thiết kế/\allowbreak{}tổng,\allowbreak{} hệ và nhật ký bộ giải,\allowbreak{} phần dư và sai số có định nghĩa rõ.

\examplebreak\subsubsection*{19. Lũy thừa}

\textbf{Input:} Trọng số w\_i (\allowbreak{}để trống: tất cả bằng 1);\allowbreak{} Cách giải bình phương tối thiểu;\allowbreak{} Quy tắc khi thiếu hạng;\allowbreak{} Điểm dự đoán x (\allowbreak{}không bắt buộc);\allowbreak{} Mức trình bày.

Bảng cặp $x,y$ là đầu vào riêng của lượt chạy; giữ ánh xạ nguồn nếu chọn một phần bảng.

\begin{enumerate}

\item Công bố N quan sát,\allowbreak{} trọng số và mô hình. Kiểm tra miền,\allowbreak{} dấu và hạng;\allowbreak{} không tự lọc dòng để ép phép biến đổi hợp lệ.

\item Viết phép biến đổi và ánh xạ tham số: z=\allowbreak{}ln|y|,\allowbreak{} u=\allowbreak{}ln x,\allowbreak{} z\ensuremath{\approx}c\_0+\allowbreak{}c\_1u;\allowbreak{} khôi phục A=\allowbreak{}s·exp(\allowbreak{}c\_0),\allowbreak{} b=\allowbreak{}c\_1.

\item Lập \ensuremath{\Phi}\_ij=\allowbreak{}\ensuremath{\varphi}\_j(\allowbreak{}x\_i),\allowbreak{} bảng các tích và các tổng;\allowbreak{} viết J=\allowbreak{}\ensuremath{\sum}w\_i(\allowbreak{}\ensuremath{\sum}c\_j\ensuremath{\varphi}\_j(\allowbreak{}x\_i)\ensuremath{-}z\_i)².

\item Từ \ensuremath{\partial}J/\allowbreak{}\ensuremath{\partial}c\_j=\allowbreak{}0 suy ra (\allowbreak{}\ensuremath{\Phi}\ensuremath{{}^T}W\ensuremath{\Phi})c=\allowbreak{}\ensuremath{\Phi}\ensuremath{{}^T}Wz. Trình bày hệ tổng quát rồi hệ đã thế số.

\item Tính nghiệm bằng bộ giải đã chọn: Gauss trên hệ chuẩn,\allowbreak{} QR trên thiết kế có trọng số hoặc SVD. Nhật ký phải đúng bộ giải.

\item Khôi phục tham số và hàm g theo biến gốc;\allowbreak{} lập r\_i=\allowbreak{}y\_i\ensuremath{-}g(\allowbreak{}x\_i),\allowbreak{} tính SSE,\allowbreak{} RMSE và mục tiêu trên thang khớp riêng biệt.

\end{enumerate}

\textbf{Output:} Hàm g(\allowbreak{}x),\allowbreak{} tham số gốc,\allowbreak{} bảng thiết kế/\allowbreak{}tổng,\allowbreak{} hệ và nhật ký bộ giải,\allowbreak{} phần dư và sai số có định nghĩa rõ.

\examplebreak\subsubsection*{20. Hàm mũ}

\textbf{Input:} Các hàm cơ sở (\allowbreak{}ngăn bằng dấu ;\allowbreak{});\allowbreak{} Trọng số w\_i (\allowbreak{}để trống: tất cả bằng 1);\allowbreak{} Cách giải bình phương tối thiểu;\allowbreak{} Quy tắc khi thiếu hạng;\allowbreak{} Điểm dự đoán x (\allowbreak{}không bắt buộc);\allowbreak{} Mức trình bày.

Bảng cặp $x,y$ là đầu vào riêng của lượt chạy; giữ ánh xạ nguồn nếu chọn một phần bảng.

\begin{enumerate}

\item Công bố N quan sát,\allowbreak{} trọng số và mô hình. Kiểm tra miền,\allowbreak{} dấu và hạng;\allowbreak{} không tự lọc dòng để ép phép biến đổi hợp lệ.

\item Viết phép biến đổi và ánh xạ tham số: z=\allowbreak{}ln|y|,\allowbreak{} z\ensuremath{\approx}c\_0+\allowbreak{}c\_1\ensuremath{\varphi}\_1(\allowbreak{}x)+\allowbreak{}…;\allowbreak{} khôi phục A=\allowbreak{}s·exp(\allowbreak{}c\_0).

\item Lập \ensuremath{\Phi}\_ij=\allowbreak{}\ensuremath{\varphi}\_j(\allowbreak{}x\_i),\allowbreak{} bảng các tích và các tổng;\allowbreak{} viết J=\allowbreak{}\ensuremath{\sum}w\_i(\allowbreak{}\ensuremath{\sum}c\_j\ensuremath{\varphi}\_j(\allowbreak{}x\_i)\ensuremath{-}z\_i)².

\item Từ \ensuremath{\partial}J/\allowbreak{}\ensuremath{\partial}c\_j=\allowbreak{}0 suy ra (\allowbreak{}\ensuremath{\Phi}\ensuremath{{}^T}W\ensuremath{\Phi})c=\allowbreak{}\ensuremath{\Phi}\ensuremath{{}^T}Wz. Trình bày hệ tổng quát rồi hệ đã thế số.

\item Tính nghiệm bằng bộ giải đã chọn: Gauss trên hệ chuẩn,\allowbreak{} QR trên thiết kế có trọng số hoặc SVD. Nhật ký phải đúng bộ giải.

\item Khôi phục tham số và hàm g theo biến gốc;\allowbreak{} lập r\_i=\allowbreak{}y\_i\ensuremath{-}g(\allowbreak{}x\_i),\allowbreak{} tính SSE,\allowbreak{} RMSE và mục tiêu trên thang khớp riêng biệt.

\end{enumerate}

\textbf{Output:} Hàm g(\allowbreak{}x),\allowbreak{} tham số gốc,\allowbreak{} bảng thiết kế/\allowbreak{}tổng,\allowbreak{} hệ và nhật ký bộ giải,\allowbreak{} phần dư và sai số có định nghĩa rõ.

\examplebreak\subsubsection*{21. Lũy thừa — SSE thang y}

\textbf{Input:} Tham số đầu A b1 b2... (\allowbreak{}để trống: từ log);\allowbreak{} Ngưỡng dừng phi tuyến;\allowbreak{} Số vòng tối đa;\allowbreak{} Trọng số w\_i (\allowbreak{}để trống: tất cả bằng 1);\allowbreak{} Bộ giải phi tuyến;\allowbreak{} Điểm dự đoán x (\allowbreak{}không bắt buộc);\allowbreak{} Mức trình bày.

Bảng cặp $x,y$ là đầu vào riêng của lượt chạy; giữ ánh xạ nguồn nếu chọn một phần bảng.

\begin{enumerate}

\item Công bố N quan sát,\allowbreak{} trọng số và mô hình. Kiểm tra miền,\allowbreak{} dấu và hạng;\allowbreak{} không tự lọc dòng để ép phép biến đổi hợp lệ.

\item Viết phép biến đổi và ánh xạ tham số: g(\allowbreak{}x)=\allowbreak{}A·exp(\allowbreak{}b ln x)=\allowbreak{}A x\textasciicircum{}b;\allowbreak{} tối thiểu SSE trên y bằng LM.

\item Công bố tham số đầu. Mỗi bước thử lập dự đoán,\allowbreak{} phần dư,\allowbreak{} Jacobian và hệ tăng cường LM với damping \ensuremath{\lambda}.

\item Ghi tham số cũ/\allowbreak{}đề xuất,\allowbreak{} SSE trước/\allowbreak{}sau,\allowbreak{} \ensuremath{\lambda} và quyết định nhận hoặc từ chối. Giữ cả các bước thử bị từ chối.

\item Kiểm tra điều kiện dừng thực tế;\allowbreak{} kết luận phụ thuộc điểm đầu,\allowbreak{} không suy ra cực tiểu toàn cục.

\item Khôi phục tham số và hàm g theo biến gốc;\allowbreak{} lập r\_i=\allowbreak{}y\_i\ensuremath{-}g(\allowbreak{}x\_i),\allowbreak{} tính SSE,\allowbreak{} RMSE và mục tiêu trên thang khớp riêng biệt.

\end{enumerate}

\textbf{Output:} Hàm g(\allowbreak{}x),\allowbreak{} tham số gốc,\allowbreak{} bảng thiết kế/\allowbreak{}tổng,\allowbreak{} hệ và nhật ký bộ giải,\allowbreak{} phần dư và sai số có định nghĩa rõ.

\examplebreak\subsubsection*{22. Hàm mũ — SSE thang y}

\textbf{Input:} Các hàm cơ sở (\allowbreak{}ngăn bằng dấu ;\allowbreak{});\allowbreak{} Tham số đầu A b1 b2... (\allowbreak{}để trống: từ log);\allowbreak{} Ngưỡng dừng phi tuyến;\allowbreak{} Số vòng tối đa;\allowbreak{} Trọng số w\_i (\allowbreak{}để trống: tất cả bằng 1);\allowbreak{} Bộ giải phi tuyến;\allowbreak{} Điểm dự đoán x (\allowbreak{}không bắt buộc);\allowbreak{} Mức trình bày.

Bảng cặp $x,y$ là đầu vào riêng của lượt chạy; giữ ánh xạ nguồn nếu chọn một phần bảng.

\begin{enumerate}

\item Công bố N quan sát,\allowbreak{} trọng số và mô hình. Kiểm tra miền,\allowbreak{} dấu và hạng;\allowbreak{} không tự lọc dòng để ép phép biến đổi hợp lệ.

\item Viết phép biến đổi và ánh xạ tham số: g(\allowbreak{}x)=\allowbreak{}A·exp(\allowbreak{}b\_1\ensuremath{\varphi}\_1(\allowbreak{}x)+\allowbreak{}…);\allowbreak{} tối thiểu SSE trên y bằng LM.

\item Công bố tham số đầu. Mỗi bước thử lập dự đoán,\allowbreak{} phần dư,\allowbreak{} Jacobian và hệ tăng cường LM với damping \ensuremath{\lambda}.

\item Ghi tham số cũ/\allowbreak{}đề xuất,\allowbreak{} SSE trước/\allowbreak{}sau,\allowbreak{} \ensuremath{\lambda} và quyết định nhận hoặc từ chối. Giữ cả các bước thử bị từ chối.

\item Kiểm tra điều kiện dừng thực tế;\allowbreak{} kết luận phụ thuộc điểm đầu,\allowbreak{} không suy ra cực tiểu toàn cục.

\item Khôi phục tham số và hàm g theo biến gốc;\allowbreak{} lập r\_i=\allowbreak{}y\_i\ensuremath{-}g(\allowbreak{}x\_i),\allowbreak{} tính SSE,\allowbreak{} RMSE và mục tiêu trên thang khớp riêng biệt.

\end{enumerate}

\textbf{Output:} Hàm g(\allowbreak{}x),\allowbreak{} tham số gốc,\allowbreak{} bảng thiết kế/\allowbreak{}tổng,\allowbreak{} hệ và nhật ký bộ giải,\allowbreak{} phần dư và sai số có định nghĩa rõ.

\examplebreak\subsubsection*{23. Log của tổ hợp cơ sở}

\textbf{Input:} Các hàm cơ sở (\allowbreak{}ngăn bằng dấu ;\allowbreak{});\allowbreak{} Trọng số w\_i (\allowbreak{}để trống: tất cả bằng 1);\allowbreak{} Cách giải bình phương tối thiểu;\allowbreak{} Quy tắc khi thiếu hạng;\allowbreak{} Điểm dự đoán x (\allowbreak{}không bắt buộc);\allowbreak{} Mức trình bày.

Bảng cặp $x,y$ là đầu vào riêng của lượt chạy; giữ ánh xạ nguồn nếu chọn một phần bảng.

\begin{enumerate}

\item Công bố N quan sát,\allowbreak{} trọng số và mô hình. Kiểm tra miền,\allowbreak{} dấu và hạng;\allowbreak{} không tự lọc dòng để ép phép biến đổi hợp lệ.

\item Viết phép biến đổi và ánh xạ tham số: z=\allowbreak{}exp(\allowbreak{}y),\allowbreak{} z\ensuremath{\approx}c\_0\ensuremath{\varphi}\_0(\allowbreak{}x)+\allowbreak{}…;\allowbreak{} khôi phục g(\allowbreak{}x)=\allowbreak{}ln(\allowbreak{}c\_0\ensuremath{\varphi}\_0(\allowbreak{}x)+\allowbreak{}…).

\item Lập \ensuremath{\Phi}\_ij=\allowbreak{}\ensuremath{\varphi}\_j(\allowbreak{}x\_i),\allowbreak{} bảng các tích và các tổng;\allowbreak{} viết J=\allowbreak{}\ensuremath{\sum}w\_i(\allowbreak{}\ensuremath{\sum}c\_j\ensuremath{\varphi}\_j(\allowbreak{}x\_i)\ensuremath{-}z\_i)².

\item Từ \ensuremath{\partial}J/\allowbreak{}\ensuremath{\partial}c\_j=\allowbreak{}0 suy ra (\allowbreak{}\ensuremath{\Phi}\ensuremath{{}^T}W\ensuremath{\Phi})c=\allowbreak{}\ensuremath{\Phi}\ensuremath{{}^T}Wz. Trình bày hệ tổng quát rồi hệ đã thế số.

\item Tính nghiệm bằng bộ giải đã chọn: Gauss trên hệ chuẩn,\allowbreak{} QR trên thiết kế có trọng số hoặc SVD. Nhật ký phải đúng bộ giải.

\item Khôi phục tham số và hàm g theo biến gốc;\allowbreak{} lập r\_i=\allowbreak{}y\_i\ensuremath{-}g(\allowbreak{}x\_i),\allowbreak{} tính SSE,\allowbreak{} RMSE và mục tiêu trên thang khớp riêng biệt.

\end{enumerate}

\textbf{Output:} Hàm g(\allowbreak{}x),\allowbreak{} tham số gốc,\allowbreak{} bảng thiết kế/\allowbreak{}tổng,\allowbreak{} hệ và nhật ký bộ giải,\allowbreak{} phần dư và sai số có định nghĩa rõ.

\examplebreak\subsubsection*{24. 1/\allowbreak{}(\allowbreak{}a+\allowbreak{}b*x)}

\textbf{Input:} Trọng số w\_i (\allowbreak{}để trống: tất cả bằng 1);\allowbreak{} Cách giải bình phương tối thiểu;\allowbreak{} Quy tắc khi thiếu hạng;\allowbreak{} Điểm dự đoán x (\allowbreak{}không bắt buộc);\allowbreak{} Mức trình bày.

Bảng cặp $x,y$ là đầu vào riêng của lượt chạy; giữ ánh xạ nguồn nếu chọn một phần bảng.

\begin{enumerate}

\item Công bố N quan sát,\allowbreak{} trọng số và mô hình. Kiểm tra miền,\allowbreak{} dấu và hạng;\allowbreak{} không tự lọc dòng để ép phép biến đổi hợp lệ.

\item Viết phép biến đổi và ánh xạ tham số: z=\allowbreak{}1/\allowbreak{}y,\allowbreak{} z\ensuremath{\approx}a+\allowbreak{}bx;\allowbreak{} khôi phục g(\allowbreak{}x)=\allowbreak{}1/\allowbreak{}(\allowbreak{}a+\allowbreak{}bx).

\item Lập \ensuremath{\Phi}\_ij=\allowbreak{}\ensuremath{\varphi}\_j(\allowbreak{}x\_i),\allowbreak{} bảng các tích và các tổng;\allowbreak{} viết J=\allowbreak{}\ensuremath{\sum}w\_i(\allowbreak{}\ensuremath{\sum}c\_j\ensuremath{\varphi}\_j(\allowbreak{}x\_i)\ensuremath{-}z\_i)².

\item Từ \ensuremath{\partial}J/\allowbreak{}\ensuremath{\partial}c\_j=\allowbreak{}0 suy ra (\allowbreak{}\ensuremath{\Phi}\ensuremath{{}^T}W\ensuremath{\Phi})c=\allowbreak{}\ensuremath{\Phi}\ensuremath{{}^T}Wz. Trình bày hệ tổng quát rồi hệ đã thế số.

\item Tính nghiệm bằng bộ giải đã chọn: Gauss trên hệ chuẩn,\allowbreak{} QR trên thiết kế có trọng số hoặc SVD. Nhật ký phải đúng bộ giải.

\item Khôi phục tham số và hàm g theo biến gốc;\allowbreak{} lập r\_i=\allowbreak{}y\_i\ensuremath{-}g(\allowbreak{}x\_i),\allowbreak{} tính SSE,\allowbreak{} RMSE và mục tiêu trên thang khớp riêng biệt.

\end{enumerate}

\textbf{Output:} Hàm g(\allowbreak{}x),\allowbreak{} tham số gốc,\allowbreak{} bảng thiết kế/\allowbreak{}tổng,\allowbreak{} hệ và nhật ký bộ giải,\allowbreak{} phần dư và sai số có định nghĩa rõ.

\examplebreak\subsubsection*{25. x/\allowbreak{}(\allowbreak{}a+\allowbreak{}b*x)}

\textbf{Input:} Trọng số w\_i (\allowbreak{}để trống: tất cả bằng 1);\allowbreak{} Cách giải bình phương tối thiểu;\allowbreak{} Quy tắc khi thiếu hạng;\allowbreak{} Điểm dự đoán x (\allowbreak{}không bắt buộc);\allowbreak{} Mức trình bày.

Bảng cặp $x,y$ là đầu vào riêng của lượt chạy; giữ ánh xạ nguồn nếu chọn một phần bảng.

\begin{enumerate}

\item Công bố N quan sát,\allowbreak{} trọng số và mô hình. Kiểm tra miền,\allowbreak{} dấu và hạng;\allowbreak{} không tự lọc dòng để ép phép biến đổi hợp lệ.

\item Viết phép biến đổi và ánh xạ tham số: z=\allowbreak{}1/\allowbreak{}y,\allowbreak{} u=\allowbreak{}1/\allowbreak{}x,\allowbreak{} z\ensuremath{\approx}c\_0+\allowbreak{}c\_1u;\allowbreak{} khôi phục a=\allowbreak{}c\_1,\allowbreak{} b=\allowbreak{}c\_0.

\item Lập \ensuremath{\Phi}\_ij=\allowbreak{}\ensuremath{\varphi}\_j(\allowbreak{}x\_i),\allowbreak{} bảng các tích và các tổng;\allowbreak{} viết J=\allowbreak{}\ensuremath{\sum}w\_i(\allowbreak{}\ensuremath{\sum}c\_j\ensuremath{\varphi}\_j(\allowbreak{}x\_i)\ensuremath{-}z\_i)².

\item Từ \ensuremath{\partial}J/\allowbreak{}\ensuremath{\partial}c\_j=\allowbreak{}0 suy ra (\allowbreak{}\ensuremath{\Phi}\ensuremath{{}^T}W\ensuremath{\Phi})c=\allowbreak{}\ensuremath{\Phi}\ensuremath{{}^T}Wz. Trình bày hệ tổng quát rồi hệ đã thế số.

\item Tính nghiệm bằng bộ giải đã chọn: Gauss trên hệ chuẩn,\allowbreak{} QR trên thiết kế có trọng số hoặc SVD. Nhật ký phải đúng bộ giải.

\item Khôi phục tham số và hàm g theo biến gốc;\allowbreak{} lập r\_i=\allowbreak{}y\_i\ensuremath{-}g(\allowbreak{}x\_i),\allowbreak{} tính SSE,\allowbreak{} RMSE và mục tiêu trên thang khớp riêng biệt.

\end{enumerate}

\textbf{Output:} Hàm g(\allowbreak{}x),\allowbreak{} tham số gốc,\allowbreak{} bảng thiết kế/\allowbreak{}tổng,\allowbreak{} hệ và nhật ký bộ giải,\allowbreak{} phần dư và sai số có định nghĩa rõ.

\examplebreak\subsubsection*{26. a+\allowbreak{}b*ln(\allowbreak{}x)}

\textbf{Input:} Trọng số w\_i (\allowbreak{}để trống: tất cả bằng 1);\allowbreak{} Cách giải bình phương tối thiểu;\allowbreak{} Quy tắc khi thiếu hạng;\allowbreak{} Điểm dự đoán x (\allowbreak{}không bắt buộc);\allowbreak{} Mức trình bày.

Bảng cặp $x,y$ là đầu vào riêng của lượt chạy; giữ ánh xạ nguồn nếu chọn một phần bảng.

\begin{enumerate}

\item Công bố N quan sát,\allowbreak{} trọng số và mô hình. Kiểm tra miền,\allowbreak{} dấu và hạng;\allowbreak{} không tự lọc dòng để ép phép biến đổi hợp lệ.

\item Viết phép biến đổi và ánh xạ tham số: z=\allowbreak{}y,\allowbreak{} u=\allowbreak{}ln x,\allowbreak{} z\ensuremath{\approx}a+\allowbreak{}bu;\allowbreak{} không đổi thang tung độ.

\item Lập \ensuremath{\Phi}\_ij=\allowbreak{}\ensuremath{\varphi}\_j(\allowbreak{}x\_i),\allowbreak{} bảng các tích và các tổng;\allowbreak{} viết J=\allowbreak{}\ensuremath{\sum}w\_i(\allowbreak{}\ensuremath{\sum}c\_j\ensuremath{\varphi}\_j(\allowbreak{}x\_i)\ensuremath{-}z\_i)².

\item Từ \ensuremath{\partial}J/\allowbreak{}\ensuremath{\partial}c\_j=\allowbreak{}0 suy ra (\allowbreak{}\ensuremath{\Phi}\ensuremath{{}^T}W\ensuremath{\Phi})c=\allowbreak{}\ensuremath{\Phi}\ensuremath{{}^T}Wz. Trình bày hệ tổng quát rồi hệ đã thế số.

\item Tính nghiệm bằng bộ giải đã chọn: Gauss trên hệ chuẩn,\allowbreak{} QR trên thiết kế có trọng số hoặc SVD. Nhật ký phải đúng bộ giải.

\item Khôi phục tham số và hàm g theo biến gốc;\allowbreak{} lập r\_i=\allowbreak{}y\_i\ensuremath{-}g(\allowbreak{}x\_i),\allowbreak{} tính SSE,\allowbreak{} RMSE và mục tiêu trên thang khớp riêng biệt.

\end{enumerate}

\textbf{Output:} Hàm g(\allowbreak{}x),\allowbreak{} tham số gốc,\allowbreak{} bảng thiết kế/\allowbreak{}tổng,\allowbreak{} hệ và nhật ký bộ giải,\allowbreak{} phần dư và sai số có định nghĩa rõ.

\examplebreak\subsubsection*{27. A*b\textasciicircum{}x}

\textbf{Input:} Trọng số w\_i (\allowbreak{}để trống: tất cả bằng 1);\allowbreak{} Cách giải bình phương tối thiểu;\allowbreak{} Quy tắc khi thiếu hạng;\allowbreak{} Điểm dự đoán x (\allowbreak{}không bắt buộc);\allowbreak{} Mức trình bày.

Bảng cặp $x,y$ là đầu vào riêng của lượt chạy; giữ ánh xạ nguồn nếu chọn một phần bảng.

\begin{enumerate}

\item Công bố N quan sát,\allowbreak{} trọng số và mô hình. Kiểm tra miền,\allowbreak{} dấu và hạng;\allowbreak{} không tự lọc dòng để ép phép biến đổi hợp lệ.

\item Viết phép biến đổi và ánh xạ tham số: z=\allowbreak{}ln|y|,\allowbreak{} z\ensuremath{\approx}c\_0+\allowbreak{}c\_1x;\allowbreak{} khôi phục A=\allowbreak{}s·exp(\allowbreak{}c\_0),\allowbreak{} B=\allowbreak{}exp(\allowbreak{}c\_1)>0.

\item Lập \ensuremath{\Phi}\_ij=\allowbreak{}\ensuremath{\varphi}\_j(\allowbreak{}x\_i),\allowbreak{} bảng các tích và các tổng;\allowbreak{} viết J=\allowbreak{}\ensuremath{\sum}w\_i(\allowbreak{}\ensuremath{\sum}c\_j\ensuremath{\varphi}\_j(\allowbreak{}x\_i)\ensuremath{-}z\_i)².

\item Từ \ensuremath{\partial}J/\allowbreak{}\ensuremath{\partial}c\_j=\allowbreak{}0 suy ra (\allowbreak{}\ensuremath{\Phi}\ensuremath{{}^T}W\ensuremath{\Phi})c=\allowbreak{}\ensuremath{\Phi}\ensuremath{{}^T}Wz. Trình bày hệ tổng quát rồi hệ đã thế số.

\item Tính nghiệm bằng bộ giải đã chọn: Gauss trên hệ chuẩn,\allowbreak{} QR trên thiết kế có trọng số hoặc SVD. Nhật ký phải đúng bộ giải.

\item Khôi phục tham số và hàm g theo biến gốc;\allowbreak{} lập r\_i=\allowbreak{}y\_i\ensuremath{-}g(\allowbreak{}x\_i),\allowbreak{} tính SSE,\allowbreak{} RMSE và mục tiêu trên thang khớp riêng biệt.

\end{enumerate}

\textbf{Output:} Hàm g(\allowbreak{}x),\allowbreak{} tham số gốc,\allowbreak{} bảng thiết kế/\allowbreak{}tổng,\allowbreak{} hệ và nhật ký bộ giải,\allowbreak{} phần dư và sai số có định nghĩa rõ.

\subsection{Nội suy ngược}

\examplebreak\subsubsection*{28. Lặp tiến}

\textbf{Input:} Giá trị y cần tìm ngược;\allowbreak{} Cận x trái (\allowbreak{}trống: khoảng đổi dấu duy nhất);\allowbreak{} Cận x phải (\allowbreak{}trống: khoảng đổi dấu duy nhất);\allowbreak{} Dung sai vị trí/\allowbreak{}bước theo x;\allowbreak{} Dung sai phần dư;\allowbreak{} Mốc dùng (\allowbreak{}JSON,\allowbreak{} trống: toàn bộ bảng);\allowbreak{} Số vòng tối đa;\allowbreak{} Quy tắc khởi tạo;\allowbreak{} t₀ khi khởi tạo thủ công;\allowbreak{} Mức trình bày.

Bảng cặp $x,y$ là đầu vào riêng của lượt chạy; giữ ánh xạ nguồn nếu chọn một phần bảng.

\begin{enumerate}

\item Công bố mức y*,\allowbreak{} các mốc được chọn và nhánh/\allowbreak{}khoảng nghiệm;\allowbreak{} kiểm tra trường hợp trùng mốc,\allowbreak{} nghiệm nhiều nhánh và dữ liệu phẳng.

\item Giữ toàn bộ tập mốc đã chọn;\allowbreak{} dựng đa thức thuận p. Lặp tiến lấy gốc đầu tập,\allowbreak{} lặp lùi lấy gốc cuối tập;\allowbreak{} khoảng nghiệm có thể nằm giữa.

\item Đặt t=\allowbreak{}(\allowbreak{}x\ensuremath{-}x\_gốc)/\allowbreak{}h,\allowbreak{} t\_0 theo quy tắc khởi tạo;\allowbreak{} tách hạng tuyến tính để viết t\_(\allowbreak{}k+\allowbreak{}1)=\allowbreak{}H(\allowbreak{}t\_k),\allowbreak{} cố định các sai phân.

\item Mỗi vòng ghi t\_k,\allowbreak{} từng tích và đóng góp phi tuyến,\allowbreak{} H(\allowbreak{}t\_k),\allowbreak{} x\_(\allowbreak{}k+\allowbreak{}1),\allowbreak{} độ đổi theo x và phần dư |p(\allowbreak{}x)\ensuremath{-}y*|.

\item Dừng thực hành cần đạt cả độ đổi theo x và phần dư. Chế độ có cận co còn kiểm tra miền bất biến,\allowbreak{} q<1 và cận sai số;\allowbreak{} giữ nhật ký thất bại.

\end{enumerate}

\textbf{Output:} Nghiệm x xấp xỉ hoặc lý do chưa đạt,\allowbreak{} tập mốc/\allowbreak{}gốc thực dùng,\allowbreak{} bảng hệ số,\allowbreak{} toàn bộ vòng lặp và tiêu chí dừng.

\examplebreak\subsubsection*{29. Lặp lùi}

\textbf{Input:} Giá trị y cần tìm ngược;\allowbreak{} Cận x trái (\allowbreak{}trống: khoảng đổi dấu duy nhất);\allowbreak{} Cận x phải (\allowbreak{}trống: khoảng đổi dấu duy nhất);\allowbreak{} Dung sai vị trí/\allowbreak{}bước theo x;\allowbreak{} Dung sai phần dư;\allowbreak{} Mốc dùng (\allowbreak{}JSON,\allowbreak{} trống: toàn bộ bảng);\allowbreak{} Số vòng tối đa;\allowbreak{} Quy tắc khởi tạo;\allowbreak{} t₀ khi khởi tạo thủ công;\allowbreak{} Mức trình bày.

Bảng cặp $x,y$ là đầu vào riêng của lượt chạy; giữ ánh xạ nguồn nếu chọn một phần bảng.

\begin{enumerate}

\item Công bố mức y*,\allowbreak{} các mốc được chọn và nhánh/\allowbreak{}khoảng nghiệm;\allowbreak{} kiểm tra trường hợp trùng mốc,\allowbreak{} nghiệm nhiều nhánh và dữ liệu phẳng.

\item Giữ toàn bộ tập mốc đã chọn;\allowbreak{} dựng đa thức thuận p. Lặp tiến lấy gốc đầu tập,\allowbreak{} lặp lùi lấy gốc cuối tập;\allowbreak{} khoảng nghiệm có thể nằm giữa.

\item Đặt t=\allowbreak{}(\allowbreak{}x\ensuremath{-}x\_gốc)/\allowbreak{}h,\allowbreak{} t\_0 theo quy tắc khởi tạo;\allowbreak{} tách hạng tuyến tính để viết t\_(\allowbreak{}k+\allowbreak{}1)=\allowbreak{}H(\allowbreak{}t\_k),\allowbreak{} cố định các sai phân.

\item Mỗi vòng ghi t\_k,\allowbreak{} từng tích và đóng góp phi tuyến,\allowbreak{} H(\allowbreak{}t\_k),\allowbreak{} x\_(\allowbreak{}k+\allowbreak{}1),\allowbreak{} độ đổi theo x và phần dư |p(\allowbreak{}x)\ensuremath{-}y*|.

\item Dừng thực hành cần đạt cả độ đổi theo x và phần dư. Chế độ có cận co còn kiểm tra miền bất biến,\allowbreak{} q<1 và cận sai số;\allowbreak{} giữ nhật ký thất bại.

\end{enumerate}

\textbf{Output:} Nghiệm x xấp xỉ hoặc lý do chưa đạt,\allowbreak{} tập mốc/\allowbreak{}gốc thực dùng,\allowbreak{} bảng hệ số,\allowbreak{} toàn bộ vòng lặp và tiêu chí dừng.

\examplebreak\subsubsection*{30. Lặp có cận co}

\textbf{Input:} Giá trị y cần tìm ngược;\allowbreak{} Cận x trái của khoảng nghiệm;\allowbreak{} Cận x phải của khoảng nghiệm;\allowbreak{} Dung sai vị trí;\allowbreak{} Dung sai phần dư;\allowbreak{} Hướng nội suy;\allowbreak{} Mốc dùng (\allowbreak{}JSON,\allowbreak{} trống: toàn bộ bảng);\allowbreak{} Số vòng tối đa;\allowbreak{} Quy tắc khởi tạo;\allowbreak{} t₀ khi khởi tạo thủ công;\allowbreak{} Mức trình bày.

Bảng cặp $x,y$ là đầu vào riêng của lượt chạy; giữ ánh xạ nguồn nếu chọn một phần bảng.

\begin{enumerate}

\item Công bố mức y*,\allowbreak{} các mốc được chọn và nhánh/\allowbreak{}khoảng nghiệm;\allowbreak{} kiểm tra trường hợp trùng mốc,\allowbreak{} nghiệm nhiều nhánh và dữ liệu phẳng.

\item Giữ toàn bộ tập mốc đã chọn;\allowbreak{} dựng đa thức thuận p. Lặp tiến lấy gốc đầu tập,\allowbreak{} lặp lùi lấy gốc cuối tập;\allowbreak{} khoảng nghiệm có thể nằm giữa.

\item Đặt t=\allowbreak{}(\allowbreak{}x\ensuremath{-}x\_gốc)/\allowbreak{}h,\allowbreak{} t\_0 theo quy tắc khởi tạo;\allowbreak{} tách hạng tuyến tính để viết t\_(\allowbreak{}k+\allowbreak{}1)=\allowbreak{}H(\allowbreak{}t\_k),\allowbreak{} cố định các sai phân.

\item Mỗi vòng ghi t\_k,\allowbreak{} từng tích và đóng góp phi tuyến,\allowbreak{} H(\allowbreak{}t\_k),\allowbreak{} x\_(\allowbreak{}k+\allowbreak{}1),\allowbreak{} độ đổi theo x và phần dư |p(\allowbreak{}x)\ensuremath{-}y*|.

\item Dừng thực hành cần đạt cả độ đổi theo x và phần dư. Chế độ có cận co còn kiểm tra miền bất biến,\allowbreak{} q<1 và cận sai số;\allowbreak{} giữ nhật ký thất bại.

\end{enumerate}

\textbf{Output:} Nghiệm x xấp xỉ hoặc lý do chưa đạt,\allowbreak{} tập mốc/\allowbreak{}gốc thực dùng,\allowbreak{} bảng hệ số,\allowbreak{} toàn bộ vòng lặp và tiêu chí dừng.

\examplebreak\subsubsection*{31. Hàm ngược — Newton}

\textbf{Input:} Giá trị y cần tìm ngược;\allowbreak{} Số mốc đã chọn (\allowbreak{}để trống: toàn tập nhập);\allowbreak{} Mốc dùng (\allowbreak{}JSON,\allowbreak{} trống: toàn bộ bảng);\allowbreak{} Mức trình bày.

Bảng cặp $x,y$ là đầu vào riêng của lượt chạy; giữ ánh xạ nguồn nếu chọn một phần bảng.

\begin{enumerate}

\item Công bố mức y*,\allowbreak{} các mốc được chọn và nhánh/\allowbreak{}khoảng nghiệm;\allowbreak{} kiểm tra trường hợp trùng mốc,\allowbreak{} nghiệm nhiều nhánh và dữ liệu phẳng.

\item Đổi vai x và y trên nhánh đã chọn;\allowbreak{} lập đa thức x\ensuremath{\approx}Q(\allowbreak{}y) bằng Newton.

\item Thế y* vào Q;\allowbreak{} trình bày các tích/\allowbreak{}hệ số và tổng từng hạng. Không gọi đây là nghiệm chính xác của p(\allowbreak{}x)=\allowbreak{}y*.

\item Kiểm tra miền nhánh và báo ngoại suy nếu mức đích nằm ngoài dữ liệu được chọn.

\end{enumerate}

\textbf{Output:} Nghiệm x xấp xỉ hoặc lý do chưa đạt,\allowbreak{} tập mốc/\allowbreak{}gốc thực dùng,\allowbreak{} bảng hệ số,\allowbreak{} toàn bộ vòng lặp và tiêu chí dừng.

\examplebreak\subsubsection*{32. Hàm ngược — Lagrange}

\textbf{Input:} Giá trị y cần tìm ngược;\allowbreak{} Số mốc đã chọn (\allowbreak{}để trống: toàn tập nhập);\allowbreak{} Mốc dùng (\allowbreak{}JSON,\allowbreak{} trống: toàn bộ bảng);\allowbreak{} Mức trình bày.

Bảng cặp $x,y$ là đầu vào riêng của lượt chạy; giữ ánh xạ nguồn nếu chọn một phần bảng.

\begin{enumerate}

\item Công bố mức y*,\allowbreak{} các mốc được chọn và nhánh/\allowbreak{}khoảng nghiệm;\allowbreak{} kiểm tra trường hợp trùng mốc,\allowbreak{} nghiệm nhiều nhánh và dữ liệu phẳng.

\item Đổi vai x và y trên nhánh đã chọn;\allowbreak{} lập đa thức x\ensuremath{\approx}Q(\allowbreak{}y) bằng Lagrange.

\item Thế y* vào Q;\allowbreak{} trình bày các tích/\allowbreak{}hệ số và tổng từng hạng. Không gọi đây là nghiệm chính xác của p(\allowbreak{}x)=\allowbreak{}y*.

\item Kiểm tra miền nhánh và báo ngoại suy nếu mức đích nằm ngoài dữ liệu được chọn.

\end{enumerate}

\textbf{Output:} Nghiệm x xấp xỉ hoặc lý do chưa đạt,\allowbreak{} tập mốc/\allowbreak{}gốc thực dùng,\allowbreak{} bảng hệ số,\allowbreak{} toàn bộ vòng lặp và tiêu chí dừng.

\examplebreak\subsubsection*{33. Chia đôi đa thức}

\textbf{Input:} Giá trị y cần tìm ngược;\allowbreak{} Cận trái;\allowbreak{} Cận phải;\allowbreak{} Dung sai vị trí;\allowbreak{} Mốc dùng (\allowbreak{}JSON,\allowbreak{} trống: toàn bộ bảng);\allowbreak{} Mức trình bày.

Bảng cặp $x,y$ là đầu vào riêng của lượt chạy; giữ ánh xạ nguồn nếu chọn một phần bảng.

\begin{enumerate}

\item Công bố mức y*,\allowbreak{} các mốc được chọn và nhánh/\allowbreak{}khoảng nghiệm;\allowbreak{} kiểm tra trường hợp trùng mốc,\allowbreak{} nghiệm nhiều nhánh và dữ liệu phẳng.

\item Giữ toàn bộ tập mốc đã chọn;\allowbreak{} dựng đa thức thuận p. Lặp tiến lấy gốc đầu tập,\allowbreak{} lặp lùi lấy gốc cuối tập;\allowbreak{} khoảng nghiệm có thể nằm giữa.

\item Đặt F(\allowbreak{}x)=\allowbreak{}p(\allowbreak{}x)\ensuremath{-}y*;\allowbreak{} kiểm tra dấu tại hai đầu. Mỗi vòng tính c=\allowbreak{}(\allowbreak{}a+\allowbreak{}b)/\allowbreak{}2,\allowbreak{} F(\allowbreak{}c) và giữ nửa đoạn còn đổi dấu.

\item Ghi a,\allowbreak{} b,\allowbreak{} c,\allowbreak{} phần dư và độ rộng từng vòng;\allowbreak{} trả về kết quả khi đạt điều kiện dừng hoặc trạng thái chưa đạt.

\end{enumerate}

\textbf{Output:} Nghiệm x xấp xỉ hoặc lý do chưa đạt,\allowbreak{} tập mốc/\allowbreak{}gốc thực dùng,\allowbreak{} bảng hệ số,\allowbreak{} toàn bộ vòng lặp và tiêu chí dừng.

\subsection{Chọn mốc}

\examplebreak\subsubsection*{34. Mốc cách đều}

\textbf{Input:} Điểm x;\allowbreak{} Số mốc cần lấy.

Bảng cặp $x,y$ là đầu vào riêng của lượt chạy; giữ ánh xạ nguồn nếu chọn một phần bảng.

\begin{enumerate}

\item Đọc cặp x,\allowbreak{}y,\allowbreak{} giữ ID nguồn;\allowbreak{} kiểm tra cặp hợp lệ và yêu cầu điểm/\allowbreak{}số mốc.

\item Trên lưới đều chọn một cửa sổ liên tiếp gần điểm tính;\allowbreak{} công bố các mốc và kiểm tra đủ số lượng.

\item Xuất các mốc đã chọn và ID nguồn;\allowbreak{} dùng tập này làm đầu vào riêng cho phương pháp nội suy được quyết định sau đó.

\end{enumerate}

\textbf{Output:} Các mốc/\allowbreak{}nhánh,\allowbreak{} thứ tự và ID nguồn;\allowbreak{} bước chọn mốc chưa phải lời giải nội suy.

\examplebreak\subsubsection*{35. Mốc nội suy ngược}

\textbf{Input:} Giá trị y.

Bảng cặp $x,y$ là đầu vào riêng của lượt chạy; giữ ánh xạ nguồn nếu chọn một phần bảng.

\begin{enumerate}

\item Đọc cặp x,\allowbreak{}y,\allowbreak{} giữ ID nguồn;\allowbreak{} kiểm tra cặp hợp lệ và yêu cầu điểm/\allowbreak{}số mốc.

\item Chia các nhánh đơn điệu và tìm các khoảng chứa mức y*;\allowbreak{} công bố mọi lựa chọn,\allowbreak{} không tự gộp nhánh.

\item Xuất các mốc đã chọn và ID nguồn;\allowbreak{} dùng tập này làm đầu vào riêng cho phương pháp nội suy được quyết định sau đó.

\end{enumerate}

\textbf{Output:} Các mốc/\allowbreak{}nhánh,\allowbreak{} thứ tự và ID nguồn;\allowbreak{} bước chọn mốc chưa phải lời giải nội suy.

\examplebreak\subsubsection*{36. Chebyshev}

\textbf{Input:} Bậc đa thức n (\allowbreak{}sinh n+\allowbreak{}1 mốc);\allowbreak{} Đầu đoạn a;\allowbreak{} Cuối đoạn b.

\begin{enumerate}

\item Cho bậc n và đoạn [a,\allowbreak{}b];\allowbreak{} cần sinh N=\allowbreak{}n+\allowbreak{}1 mốc.

\item Tính \ensuremath{\theta}\_k=\allowbreak{}(\allowbreak{}2k+\allowbreak{}1)\ensuremath{\pi}/\allowbreak{}(\allowbreak{}2N),\allowbreak{} k=\allowbreak{}0,\allowbreak{}…,\allowbreak{}N\ensuremath{-}1;\allowbreak{} đặt x\_k=\allowbreak{}(\allowbreak{}a+\allowbreak{}b)/\allowbreak{}2+\allowbreak{}(\allowbreak{}b\ensuremath{-}a)cos \ensuremath{\theta}\_k/\allowbreak{}2.

\item Sắp mốc tăng và giữ ánh xạ với chỉ số k ban đầu;\allowbreak{} đánh giá hàm tại các mốc này nếu cần bảng y.

\end{enumerate}

\textbf{Output:} Các mốc/\allowbreak{}nhánh,\allowbreak{} thứ tự và ID nguồn;\allowbreak{} bước chọn mốc chưa phải lời giải nội suy.

\subsection{Đạo hàm \& Horner}

\examplebreak\subsubsection*{37. Đạo hàm từ bảng}

\textbf{Input:} Điểm x;\allowbreak{} Bậc p (\allowbreak{}dùng p+\allowbreak{}1 mốc);\allowbreak{} Cấp đạo hàm.

Bảng cặp $x,y$ là đầu vào riêng của lượt chạy; giữ ánh xạ nguồn nếu chọn một phần bảng.

\begin{enumerate}

\item Chọn p+\allowbreak{}1 mốc gần điểm tính,\allowbreak{} giữ ID nguồn;\allowbreak{} xác định cấp đạo hàm r và đa thức bậc không quá p.

\item Đặt tọa độ chuẩn hóa quanh điểm tính,\allowbreak{} lập các trọng số đạo hàm từ điều kiện chính xác trên đơn thức.

\item Lập bảng w\_i,\allowbreak{} y\_i,\allowbreak{} w\_i y\_i;\allowbreak{} cộng đủ mọi hạng và đổi đơn vị đạo hàm nếu có chuẩn hóa.

\item Kiểm tra tính chính xác trên các đơn thức và công bố mốc thực dùng;\allowbreak{} kết luận đạo hàm của mô hình,\allowbreak{} không tự nhận cận sai số hàm thật.

\end{enumerate}

\textbf{Output:} Giá trị/\allowbreak{}đạo hàm,\allowbreak{} bảng nhân–cộng hoặc trọng số và đối chiếu;\allowbreak{} đạo hàm vượt bậc đa thức bằng 0.

\examplebreak\subsubsection*{38. Horner}

\textbf{Input:} Hệ số theo bậc giảm dần;\allowbreak{} Điểm x;\allowbreak{} Đạo hàm đến cấp.

\begin{enumerate}

\item Đọc hệ số theo bậc giảm,\allowbreak{} giữ các hệ số 0 giữa/\allowbreak{}cuối;\allowbreak{} viết lại P(\allowbreak{}x) và điểm c.

\item Lập bảng ngang: hàng A\_j;\allowbreak{} hàng c·b\_(\allowbreak{}j\ensuremath{-}1) dịch sang phải một ô;\allowbreak{} hàng cộng b\_j=\allowbreak{}A\_j+\allowbreak{}c·b\_(\allowbreak{}j\ensuremath{-}1).

\item Phần tử cuối là số dư s\_0=\allowbreak{}P(\allowbreak{}c);\allowbreak{} các phần tử trước tạo thương Q\_1.

\item Chia tiếp thương cho x\ensuremath{-}c,\allowbreak{} lấy s\_r;\allowbreak{} P\textasciicircum{}(\allowbreak{}r)(\allowbreak{}c)=\allowbreak{}r!s\_r. Công bố số dư và thương của từng lượt.

\item Kết luận giá trị và các đạo hàm;\allowbreak{} đây là tính hữu hạn,\allowbreak{} không có điều kiện hội tụ.

\end{enumerate}

\textbf{Output:} Giá trị/\allowbreak{}đạo hàm,\allowbreak{} bảng nhân–cộng hoặc trọng số và đối chiếu;\allowbreak{} đạo hàm vượt bậc đa thức bằng 0.

\subsection{Tích phân}

\examplebreak\subsubsection*{39. Hình thang}

\textbf{Input:} Hàm f(\allowbreak{}x);\allowbreak{} Cận trái a;\allowbreak{} Cận phải b;\allowbreak{} Số khoảng chia n;\allowbreak{} Dung sai tự chia lưới (\allowbreak{}để trống: dùng n);\allowbreak{} Cận đạo hàm M (\allowbreak{}không bắt buộc);\allowbreak{} Nguồn dữ liệu.

\begin{enumerate}

\item Xác định nguồn Hàm/\allowbreak{}Bảng,\allowbreak{} hai cận,\allowbreak{} n và điều kiện miền;\allowbreak{} nếu đảo cận thì đổi dấu kết quả.

\item Lập lưới đúng quy tắc và bảng x\_i,\allowbreak{} y\_i,\allowbreak{} w\_i;\allowbreak{} h=\allowbreak{}(\allowbreak{}b\ensuremath{-}a)/\allowbreak{}n. Với trung điểm,\allowbreak{} dữ liệu phải ở tâm mỗi ô.

\item Viết công thức và thế các hạng: I\_h\ensuremath{\approx}h[y\_0+\allowbreak{}2y\_1+\allowbreak{}…+\allowbreak{}2y\_(\allowbreak{}n\ensuremath{-}1)+\allowbreak{}y\_n]/\allowbreak{}2.

\item Nếu nhập \ensuremath{\varepsilon} và M,\allowbreak{} chọn n từ cận lý thuyết. Nếu chỉ nhập \ensuremath{\varepsilon},\allowbreak{} tăng lưới và kiểm tra Runge cùng lưới đối chiếu.

\item Công bố kết quả ở lưới được trả về,\allowbreak{} các lượt chia lưới và sai số;\allowbreak{} Runge là ước lượng,\allowbreak{} cận M có giả thiết riêng.

\end{enumerate}

\textbf{Output:} Tích phân xấp xỉ,\allowbreak{} lưới/\allowbreak{}trọng số,\allowbreak{} các lượt tinh chỉnh và ước lượng/\allowbreak{}cận sai số được phân biệt rõ.

\examplebreak\subsubsection*{40. Simpson}

\textbf{Input:} Hàm f(\allowbreak{}x);\allowbreak{} Cận trái a;\allowbreak{} Cận phải b;\allowbreak{} Số khoảng chia n;\allowbreak{} Dung sai tự chia lưới (\allowbreak{}để trống: dùng n);\allowbreak{} Cận đạo hàm M (\allowbreak{}không bắt buộc);\allowbreak{} Nguồn dữ liệu.

\begin{enumerate}

\item Xác định nguồn Hàm/\allowbreak{}Bảng,\allowbreak{} hai cận,\allowbreak{} n và điều kiện miền;\allowbreak{} nếu đảo cận thì đổi dấu kết quả.

\item Lập lưới đúng quy tắc và bảng x\_i,\allowbreak{} y\_i,\allowbreak{} w\_i;\allowbreak{} h=\allowbreak{}(\allowbreak{}b\ensuremath{-}a)/\allowbreak{}n. Với trung điểm,\allowbreak{} dữ liệu phải ở tâm mỗi ô.

\item Viết công thức và thế các hạng: I\_h\ensuremath{\approx}h[y\_0+\allowbreak{}4y\_1+\allowbreak{}2y\_2+\allowbreak{}4y\_3+\allowbreak{}…+\allowbreak{}4y\_(\allowbreak{}n\ensuremath{-}1)+\allowbreak{}y\_n]/\allowbreak{}3,\allowbreak{} n chẵn.

\item Nếu nhập \ensuremath{\varepsilon} và M,\allowbreak{} chọn n từ cận lý thuyết. Nếu chỉ nhập \ensuremath{\varepsilon},\allowbreak{} tăng lưới và kiểm tra Runge cùng lưới đối chiếu.

\item Công bố kết quả ở lưới được trả về,\allowbreak{} các lượt chia lưới và sai số;\allowbreak{} Runge là ước lượng,\allowbreak{} cận M có giả thiết riêng.

\end{enumerate}

\textbf{Output:} Tích phân xấp xỉ,\allowbreak{} lưới/\allowbreak{}trọng số,\allowbreak{} các lượt tinh chỉnh và ước lượng/\allowbreak{}cận sai số được phân biệt rõ.

\examplebreak\subsubsection*{41. Trung điểm}

\textbf{Input:} Hàm f(\allowbreak{}x);\allowbreak{} Cận trái a;\allowbreak{} Cận phải b;\allowbreak{} Số khoảng chia n;\allowbreak{} Dung sai tự chia lưới (\allowbreak{}để trống: dùng n);\allowbreak{} Cận đạo hàm M (\allowbreak{}không bắt buộc);\allowbreak{} Nguồn dữ liệu.

\begin{enumerate}

\item Xác định nguồn Hàm/\allowbreak{}Bảng,\allowbreak{} hai cận,\allowbreak{} n và điều kiện miền;\allowbreak{} nếu đảo cận thì đổi dấu kết quả.

\item Lập lưới đúng quy tắc và bảng x\_i,\allowbreak{} y\_i,\allowbreak{} w\_i;\allowbreak{} h=\allowbreak{}(\allowbreak{}b\ensuremath{-}a)/\allowbreak{}n. Với trung điểm,\allowbreak{} dữ liệu phải ở tâm mỗi ô.

\item Viết công thức và thế các hạng: I\_h\ensuremath{\approx}h[f(\allowbreak{}a+\allowbreak{}h/\allowbreak{}2)+\allowbreak{}f(\allowbreak{}a+\allowbreak{}3h/\allowbreak{}2)+\allowbreak{}…+\allowbreak{}f(\allowbreak{}b\ensuremath{-}h/\allowbreak{}2)].

\item Nếu nhập \ensuremath{\varepsilon} và M,\allowbreak{} chọn n từ cận lý thuyết. Nếu chỉ nhập \ensuremath{\varepsilon},\allowbreak{} tăng lưới và kiểm tra Runge cùng lưới đối chiếu.

\item Công bố kết quả ở lưới được trả về,\allowbreak{} các lượt chia lưới và sai số;\allowbreak{} Runge là ước lượng,\allowbreak{} cận M có giả thiết riêng.

\end{enumerate}

\textbf{Output:} Tích phân xấp xỉ,\allowbreak{} lưới/\allowbreak{}trọng số,\allowbreak{} các lượt tinh chỉnh và ước lượng/\allowbreak{}cận sai số được phân biệt rõ.

\examplebreak\subsubsection*{42. Newton-Cotes}

\textbf{Input:} Hàm f(\allowbreak{}x);\allowbreak{} Cận trái a;\allowbreak{} Cận phải b;\allowbreak{} Số khoảng chia n;\allowbreak{} Bậc Newton-Cotes;\allowbreak{} Dung sai tự chia lưới (\allowbreak{}để trống: dùng n);\allowbreak{} Cận đạo hàm M (\allowbreak{}không bắt buộc);\allowbreak{} Nguồn dữ liệu.

\begin{enumerate}

\item Xác định nguồn Hàm/\allowbreak{}Bảng,\allowbreak{} hai cận,\allowbreak{} n và điều kiện miền;\allowbreak{} nếu đảo cận thì đổi dấu kết quả.

\item Lập lưới đúng quy tắc và bảng x\_i,\allowbreak{} y\_i,\allowbreak{} w\_i;\allowbreak{} h=\allowbreak{}(\allowbreak{}b\ensuremath{-}a)/\allowbreak{}n. Với trung điểm,\allowbreak{} dữ liệu phải ở tâm mỗi ô.

\item Viết công thức và thế các hạng: I\_h\ensuremath{\approx}h[w\_0y\_0+\allowbreak{}w\_1y\_1+\allowbreak{}…+\allowbreak{}w\_ny\_n],\allowbreak{} n chia hết cho bậc khối.

\item Nếu nhập \ensuremath{\varepsilon} và M,\allowbreak{} chọn n từ cận lý thuyết. Nếu chỉ nhập \ensuremath{\varepsilon},\allowbreak{} tăng lưới và kiểm tra Runge cùng lưới đối chiếu.

\item Công bố kết quả ở lưới được trả về,\allowbreak{} các lượt chia lưới và sai số;\allowbreak{} Runge là ước lượng,\allowbreak{} cận M có giả thiết riêng.

\end{enumerate}

\textbf{Output:} Tích phân xấp xỉ,\allowbreak{} lưới/\allowbreak{}trọng số,\allowbreak{} các lượt tinh chỉnh và ước lượng/\allowbreak{}cận sai số được phân biệt rõ.

\examplebreak\subsubsection*{43. Tích phân kép}

\textbf{Input:} Hàm f(\allowbreak{}x,\allowbreak{}y);\allowbreak{} Cận trái a;\allowbreak{} Cận phải b;\allowbreak{} Số khoảng chia n;\allowbreak{} Cận dưới y;\allowbreak{} Cận trên y;\allowbreak{} Quy tắc;\allowbreak{} Dung sai Runge.

\begin{enumerate}

\item Công bố hình chữ nhật,\allowbreak{} hàm hai biến,\allowbreak{} quy tắc và n trên mỗi trục;\allowbreak{} kiểm tra miền hàm.

\item Lập hai lưới và trọng số một chiều;\allowbreak{} trọng số tại (\allowbreak{}i,\allowbreak{}j) là tích w\_i v\_j.

\item Cộng từng hàng rồi cộng các hàng: I\ensuremath{\approx}h\_x h\_y \ensuremath{\sum}\_i\ensuremath{\sum}\_j w\_i v\_j f(\allowbreak{}x\_i,\allowbreak{}y\_j).

\item Chia đôi bước ở cả hai trục;\allowbreak{} ghi n,\allowbreak{} I và Runge từng vòng. Đối chiếu thêm lưới không lồng nhau khi cần.

\item Kết luận xấp xỉ và ước lượng Runge;\allowbreak{} điều kiện trơn của hàm vẫn phải được xét.

\end{enumerate}

\textbf{Output:} Tích phân xấp xỉ,\allowbreak{} lưới/\allowbreak{}trọng số,\allowbreak{} các lượt tinh chỉnh và ước lượng/\allowbreak{}cận sai số được phân biệt rõ.

\subsection{Phương trình vi phân}

\examplebreak\subsubsection*{44. Euler hiện}

\textbf{Input:} Vế phải x\ensuremath{\prime};\allowbreak{} y\ensuremath{\prime};\allowbreak{} z\ensuremath{\prime} (\allowbreak{}1–3 hàm);\allowbreak{} Giá trị đầu x,\allowbreak{} y,\allowbreak{} z;\allowbreak{} Thời điểm đầu t0;\allowbreak{} Thời điểm cuối;\allowbreak{} Bước h;\allowbreak{} Dung sai bước ẩn.

\begin{enumerate}

\item Viết hệ u\ensuremath{\prime}=\allowbreak{}f(\allowbreak{}t,\allowbreak{}u),\allowbreak{} u(\allowbreak{}a)=\allowbreak{}u\_0;\allowbreak{} trong ô nhập,\allowbreak{} biến độc lập là t và các trạng thái là x,\allowbreak{}y,\allowbreak{}z.

\item Lập lưới t\_i=\allowbreak{}a+\allowbreak{}ih;\allowbreak{} h phải đúng chiều và kết thúc đúng ở b. Công bố bộ hệ số của phương pháp.

\item Công bố bảng A,\allowbreak{}b,\allowbreak{}c. Mỗi bước tính lần lượt K\_j=\allowbreak{}f(\allowbreak{}t\_i+\allowbreak{}c\_jh,\allowbreak{}u\_i+\allowbreak{}h\ensuremath{\sum}\_(\allowbreak{}l<j)a\_jl K\_l),\allowbreak{} rồi u\_(\allowbreak{}i+\allowbreak{}1)=\allowbreak{}u\_i+\allowbreak{}h\ensuremath{\sum}\_j b\_jK\_j.

\item Ghi các tầng/\allowbreak{}dự báo/\allowbreak{}hiệu chỉnh thực sự được tính. Bước ẩn kiểm tra phần dư và Jacobian;\allowbreak{} đánh giá của bộ giải không đồng nghĩa một vòng Newton.

\item Lập bảng nghiệm qua mọi t\_i,\allowbreak{} kèm phần dư bước ẩn. Kết luận trên lưới h;\allowbreak{} dung sai hệ ẩn không phải cận sai số nghiệm ODE.

\end{enumerate}

\textbf{Output:} Nghiệm tại các mốc,\allowbreak{} bảng hệ số/\allowbreak{}tầng RK hoặc dự báo–hiệu chỉnh,\allowbreak{} phần dư và thông tin kiểm tra bước ẩn.

\examplebreak\subsubsection*{45. Euler ẩn}

\textbf{Input:} Vế phải x\ensuremath{\prime};\allowbreak{} y\ensuremath{\prime};\allowbreak{} z\ensuremath{\prime} (\allowbreak{}1–3 hàm);\allowbreak{} Giá trị đầu x,\allowbreak{} y,\allowbreak{} z;\allowbreak{} Thời điểm đầu t0;\allowbreak{} Thời điểm cuối;\allowbreak{} Bước h;\allowbreak{} Dung sai bước ẩn.

\begin{enumerate}

\item Viết hệ u\ensuremath{\prime}=\allowbreak{}f(\allowbreak{}t,\allowbreak{}u),\allowbreak{} u(\allowbreak{}a)=\allowbreak{}u\_0;\allowbreak{} trong ô nhập,\allowbreak{} biến độc lập là t và các trạng thái là x,\allowbreak{}y,\allowbreak{}z.

\item Lập lưới t\_i=\allowbreak{}a+\allowbreak{}ih;\allowbreak{} h phải đúng chiều và kết thúc đúng ở b. Công bố bộ hệ số của phương pháp.

\item Giải R(\allowbreak{}v)=\allowbreak{}v\ensuremath{-}u\_i\ensuremath{-}h f(\allowbreak{}t\_(\allowbreak{}i+\allowbreak{}1),\allowbreak{}v)=\allowbreak{}0 bằng bộ giải hệ ẩn.

\item Ghi các tầng/\allowbreak{}dự báo/\allowbreak{}hiệu chỉnh thực sự được tính. Bước ẩn kiểm tra phần dư và Jacobian;\allowbreak{} đánh giá của bộ giải không đồng nghĩa một vòng Newton.

\item Lập bảng nghiệm qua mọi t\_i,\allowbreak{} kèm phần dư bước ẩn. Kết luận trên lưới h;\allowbreak{} dung sai hệ ẩn không phải cận sai số nghiệm ODE.

\end{enumerate}

\textbf{Output:} Nghiệm tại các mốc,\allowbreak{} bảng hệ số/\allowbreak{}tầng RK hoặc dự báo–hiệu chỉnh,\allowbreak{} phần dư và thông tin kiểm tra bước ẩn.

\examplebreak\subsubsection*{46. Hình thang}

\textbf{Input:} Vế phải x\ensuremath{\prime};\allowbreak{} y\ensuremath{\prime};\allowbreak{} z\ensuremath{\prime} (\allowbreak{}1–3 hàm);\allowbreak{} Giá trị đầu x,\allowbreak{} y,\allowbreak{} z;\allowbreak{} Thời điểm đầu t0;\allowbreak{} Thời điểm cuối;\allowbreak{} Bước h;\allowbreak{} Dung sai bước ẩn.

\begin{enumerate}

\item Viết hệ u\ensuremath{\prime}=\allowbreak{}f(\allowbreak{}t,\allowbreak{}u),\allowbreak{} u(\allowbreak{}a)=\allowbreak{}u\_0;\allowbreak{} trong ô nhập,\allowbreak{} biến độc lập là t và các trạng thái là x,\allowbreak{}y,\allowbreak{}z.

\item Lập lưới t\_i=\allowbreak{}a+\allowbreak{}ih;\allowbreak{} h phải đúng chiều và kết thúc đúng ở b. Công bố bộ hệ số của phương pháp.

\item Giải R(\allowbreak{}v)=\allowbreak{}v\ensuremath{-}u\_i\ensuremath{-}h[f(\allowbreak{}t\_i,\allowbreak{}u\_i)+\allowbreak{}f(\allowbreak{}t\_(\allowbreak{}i+\allowbreak{}1),\allowbreak{}v)]/\allowbreak{}2=\allowbreak{}0 bằng bộ giải hệ ẩn.

\item Ghi các tầng/\allowbreak{}dự báo/\allowbreak{}hiệu chỉnh thực sự được tính. Bước ẩn kiểm tra phần dư và Jacobian;\allowbreak{} đánh giá của bộ giải không đồng nghĩa một vòng Newton.

\item Lập bảng nghiệm qua mọi t\_i,\allowbreak{} kèm phần dư bước ẩn. Kết luận trên lưới h;\allowbreak{} dung sai hệ ẩn không phải cận sai số nghiệm ODE.

\end{enumerate}

\textbf{Output:} Nghiệm tại các mốc,\allowbreak{} bảng hệ số/\allowbreak{}tầng RK hoặc dự báo–hiệu chỉnh,\allowbreak{} phần dư và thông tin kiểm tra bước ẩn.

\examplebreak\subsubsection*{47. RK2}

\textbf{Input:} Vế phải x\ensuremath{\prime};\allowbreak{} y\ensuremath{\prime};\allowbreak{} z\ensuremath{\prime} (\allowbreak{}1–3 hàm);\allowbreak{} Giá trị đầu x,\allowbreak{} y,\allowbreak{} z;\allowbreak{} Thời điểm đầu t0;\allowbreak{} Thời điểm cuối;\allowbreak{} Bước h;\allowbreak{} Dung sai bước ẩn.

\begin{enumerate}

\item Viết hệ u\ensuremath{\prime}=\allowbreak{}f(\allowbreak{}t,\allowbreak{}u),\allowbreak{} u(\allowbreak{}a)=\allowbreak{}u\_0;\allowbreak{} trong ô nhập,\allowbreak{} biến độc lập là t và các trạng thái là x,\allowbreak{}y,\allowbreak{}z.

\item Lập lưới t\_i=\allowbreak{}a+\allowbreak{}ih;\allowbreak{} h phải đúng chiều và kết thúc đúng ở b. Công bố bộ hệ số của phương pháp.

\item Công bố bảng A,\allowbreak{}b,\allowbreak{}c. Mỗi bước tính lần lượt K\_j=\allowbreak{}f(\allowbreak{}t\_i+\allowbreak{}c\_jh,\allowbreak{}u\_i+\allowbreak{}h\ensuremath{\sum}\_(\allowbreak{}l<j)a\_jl K\_l),\allowbreak{} rồi u\_(\allowbreak{}i+\allowbreak{}1)=\allowbreak{}u\_i+\allowbreak{}h\ensuremath{\sum}\_j b\_jK\_j.

\item Ghi các tầng/\allowbreak{}dự báo/\allowbreak{}hiệu chỉnh thực sự được tính. Bước ẩn kiểm tra phần dư và Jacobian;\allowbreak{} đánh giá của bộ giải không đồng nghĩa một vòng Newton.

\item Lập bảng nghiệm qua mọi t\_i,\allowbreak{} kèm phần dư bước ẩn. Kết luận trên lưới h;\allowbreak{} dung sai hệ ẩn không phải cận sai số nghiệm ODE.

\end{enumerate}

\textbf{Output:} Nghiệm tại các mốc,\allowbreak{} bảng hệ số/\allowbreak{}tầng RK hoặc dự báo–hiệu chỉnh,\allowbreak{} phần dư và thông tin kiểm tra bước ẩn.

\examplebreak\subsubsection*{48. RK3}

\textbf{Input:} Vế phải x\ensuremath{\prime};\allowbreak{} y\ensuremath{\prime};\allowbreak{} z\ensuremath{\prime} (\allowbreak{}1–3 hàm);\allowbreak{} Giá trị đầu x,\allowbreak{} y,\allowbreak{} z;\allowbreak{} Thời điểm đầu t0;\allowbreak{} Thời điểm cuối;\allowbreak{} Bước h;\allowbreak{} Dung sai bước ẩn.

\begin{enumerate}

\item Viết hệ u\ensuremath{\prime}=\allowbreak{}f(\allowbreak{}t,\allowbreak{}u),\allowbreak{} u(\allowbreak{}a)=\allowbreak{}u\_0;\allowbreak{} trong ô nhập,\allowbreak{} biến độc lập là t và các trạng thái là x,\allowbreak{}y,\allowbreak{}z.

\item Lập lưới t\_i=\allowbreak{}a+\allowbreak{}ih;\allowbreak{} h phải đúng chiều và kết thúc đúng ở b. Công bố bộ hệ số của phương pháp.

\item Công bố bảng A,\allowbreak{}b,\allowbreak{}c. Mỗi bước tính lần lượt K\_j=\allowbreak{}f(\allowbreak{}t\_i+\allowbreak{}c\_jh,\allowbreak{}u\_i+\allowbreak{}h\ensuremath{\sum}\_(\allowbreak{}l<j)a\_jl K\_l),\allowbreak{} rồi u\_(\allowbreak{}i+\allowbreak{}1)=\allowbreak{}u\_i+\allowbreak{}h\ensuremath{\sum}\_j b\_jK\_j.

\item Ghi các tầng/\allowbreak{}dự báo/\allowbreak{}hiệu chỉnh thực sự được tính. Bước ẩn kiểm tra phần dư và Jacobian;\allowbreak{} đánh giá của bộ giải không đồng nghĩa một vòng Newton.

\item Lập bảng nghiệm qua mọi t\_i,\allowbreak{} kèm phần dư bước ẩn. Kết luận trên lưới h;\allowbreak{} dung sai hệ ẩn không phải cận sai số nghiệm ODE.

\end{enumerate}

\textbf{Output:} Nghiệm tại các mốc,\allowbreak{} bảng hệ số/\allowbreak{}tầng RK hoặc dự báo–hiệu chỉnh,\allowbreak{} phần dư và thông tin kiểm tra bước ẩn.

\examplebreak\subsubsection*{49. RK4}

\textbf{Input:} Vế phải x\ensuremath{\prime};\allowbreak{} y\ensuremath{\prime};\allowbreak{} z\ensuremath{\prime} (\allowbreak{}1–3 hàm);\allowbreak{} Giá trị đầu x,\allowbreak{} y,\allowbreak{} z;\allowbreak{} Thời điểm đầu t0;\allowbreak{} Thời điểm cuối;\allowbreak{} Bước h;\allowbreak{} Dung sai bước ẩn.

\begin{enumerate}

\item Viết hệ u\ensuremath{\prime}=\allowbreak{}f(\allowbreak{}t,\allowbreak{}u),\allowbreak{} u(\allowbreak{}a)=\allowbreak{}u\_0;\allowbreak{} trong ô nhập,\allowbreak{} biến độc lập là t và các trạng thái là x,\allowbreak{}y,\allowbreak{}z.

\item Lập lưới t\_i=\allowbreak{}a+\allowbreak{}ih;\allowbreak{} h phải đúng chiều và kết thúc đúng ở b. Công bố bộ hệ số của phương pháp.

\item Công bố bảng A,\allowbreak{}b,\allowbreak{}c. Mỗi bước tính lần lượt K\_j=\allowbreak{}f(\allowbreak{}t\_i+\allowbreak{}c\_jh,\allowbreak{}u\_i+\allowbreak{}h\ensuremath{\sum}\_(\allowbreak{}l<j)a\_jl K\_l),\allowbreak{} rồi u\_(\allowbreak{}i+\allowbreak{}1)=\allowbreak{}u\_i+\allowbreak{}h\ensuremath{\sum}\_j b\_jK\_j.

\item Ghi các tầng/\allowbreak{}dự báo/\allowbreak{}hiệu chỉnh thực sự được tính. Bước ẩn kiểm tra phần dư và Jacobian;\allowbreak{} đánh giá của bộ giải không đồng nghĩa một vòng Newton.

\item Lập bảng nghiệm qua mọi t\_i,\allowbreak{} kèm phần dư bước ẩn. Kết luận trên lưới h;\allowbreak{} dung sai hệ ẩn không phải cận sai số nghiệm ODE.

\end{enumerate}

\textbf{Output:} Nghiệm tại các mốc,\allowbreak{} bảng hệ số/\allowbreak{}tầng RK hoặc dự báo–hiệu chỉnh,\allowbreak{} phần dư và thông tin kiểm tra bước ẩn.

\examplebreak\subsubsection*{50. Adams 2}

\textbf{Input:} Vế phải x\ensuremath{\prime};\allowbreak{} y\ensuremath{\prime};\allowbreak{} z\ensuremath{\prime} (\allowbreak{}1–3 hàm);\allowbreak{} Giá trị đầu x,\allowbreak{} y,\allowbreak{} z;\allowbreak{} Thời điểm đầu t0;\allowbreak{} Thời điểm cuối;\allowbreak{} Bước h;\allowbreak{} Dung sai bước ẩn.

\begin{enumerate}

\item Viết hệ u\ensuremath{\prime}=\allowbreak{}f(\allowbreak{}t,\allowbreak{}u),\allowbreak{} u(\allowbreak{}a)=\allowbreak{}u\_0;\allowbreak{} trong ô nhập,\allowbreak{} biến độc lập là t và các trạng thái là x,\allowbreak{}y,\allowbreak{}z.

\item Lập lưới t\_i=\allowbreak{}a+\allowbreak{}ih;\allowbreak{} h phải đúng chiều và kết thúc đúng ở b. Công bố bộ hệ số của phương pháp.

\item Khởi tạo s\ensuremath{-}1 bước bằng RK4;\allowbreak{} từ đủ lịch sử,\allowbreak{} dự báo AB\_s rồi giải phương trình AM\_(\allowbreak{}s+\allowbreak{}1) với dự báo làm điểm đầu.

\item Ghi các tầng/\allowbreak{}dự báo/\allowbreak{}hiệu chỉnh thực sự được tính. Bước ẩn kiểm tra phần dư và Jacobian;\allowbreak{} đánh giá của bộ giải không đồng nghĩa một vòng Newton.

\item Lập bảng nghiệm qua mọi t\_i,\allowbreak{} kèm phần dư bước ẩn. Kết luận trên lưới h;\allowbreak{} dung sai hệ ẩn không phải cận sai số nghiệm ODE.

\end{enumerate}

\textbf{Output:} Nghiệm tại các mốc,\allowbreak{} bảng hệ số/\allowbreak{}tầng RK hoặc dự báo–hiệu chỉnh,\allowbreak{} phần dư và thông tin kiểm tra bước ẩn.

\examplebreak\subsubsection*{51. Adams 3}

\textbf{Input:} Vế phải x\ensuremath{\prime};\allowbreak{} y\ensuremath{\prime};\allowbreak{} z\ensuremath{\prime} (\allowbreak{}1–3 hàm);\allowbreak{} Giá trị đầu x,\allowbreak{} y,\allowbreak{} z;\allowbreak{} Thời điểm đầu t0;\allowbreak{} Thời điểm cuối;\allowbreak{} Bước h;\allowbreak{} Dung sai bước ẩn.

\begin{enumerate}

\item Viết hệ u\ensuremath{\prime}=\allowbreak{}f(\allowbreak{}t,\allowbreak{}u),\allowbreak{} u(\allowbreak{}a)=\allowbreak{}u\_0;\allowbreak{} trong ô nhập,\allowbreak{} biến độc lập là t và các trạng thái là x,\allowbreak{}y,\allowbreak{}z.

\item Lập lưới t\_i=\allowbreak{}a+\allowbreak{}ih;\allowbreak{} h phải đúng chiều và kết thúc đúng ở b. Công bố bộ hệ số của phương pháp.

\item Khởi tạo s\ensuremath{-}1 bước bằng RK4;\allowbreak{} từ đủ lịch sử,\allowbreak{} dự báo AB\_s rồi giải phương trình AM\_(\allowbreak{}s+\allowbreak{}1) với dự báo làm điểm đầu.

\item Ghi các tầng/\allowbreak{}dự báo/\allowbreak{}hiệu chỉnh thực sự được tính. Bước ẩn kiểm tra phần dư và Jacobian;\allowbreak{} đánh giá của bộ giải không đồng nghĩa một vòng Newton.

\item Lập bảng nghiệm qua mọi t\_i,\allowbreak{} kèm phần dư bước ẩn. Kết luận trên lưới h;\allowbreak{} dung sai hệ ẩn không phải cận sai số nghiệm ODE.

\end{enumerate}

\textbf{Output:} Nghiệm tại các mốc,\allowbreak{} bảng hệ số/\allowbreak{}tầng RK hoặc dự báo–hiệu chỉnh,\allowbreak{} phần dư và thông tin kiểm tra bước ẩn.

\examplebreak\subsubsection*{52. Adams 4}

\textbf{Input:} Vế phải x\ensuremath{\prime};\allowbreak{} y\ensuremath{\prime};\allowbreak{} z\ensuremath{\prime} (\allowbreak{}1–3 hàm);\allowbreak{} Giá trị đầu x,\allowbreak{} y,\allowbreak{} z;\allowbreak{} Thời điểm đầu t0;\allowbreak{} Thời điểm cuối;\allowbreak{} Bước h;\allowbreak{} Dung sai bước ẩn.

\begin{enumerate}

\item Viết hệ u\ensuremath{\prime}=\allowbreak{}f(\allowbreak{}t,\allowbreak{}u),\allowbreak{} u(\allowbreak{}a)=\allowbreak{}u\_0;\allowbreak{} trong ô nhập,\allowbreak{} biến độc lập là t và các trạng thái là x,\allowbreak{}y,\allowbreak{}z.

\item Lập lưới t\_i=\allowbreak{}a+\allowbreak{}ih;\allowbreak{} h phải đúng chiều và kết thúc đúng ở b. Công bố bộ hệ số của phương pháp.

\item Khởi tạo s\ensuremath{-}1 bước bằng RK4;\allowbreak{} từ đủ lịch sử,\allowbreak{} dự báo AB\_s rồi giải phương trình AM\_(\allowbreak{}s+\allowbreak{}1) với dự báo làm điểm đầu.

\item Ghi các tầng/\allowbreak{}dự báo/\allowbreak{}hiệu chỉnh thực sự được tính. Bước ẩn kiểm tra phần dư và Jacobian;\allowbreak{} đánh giá của bộ giải không đồng nghĩa một vòng Newton.

\item Lập bảng nghiệm qua mọi t\_i,\allowbreak{} kèm phần dư bước ẩn. Kết luận trên lưới h;\allowbreak{} dung sai hệ ẩn không phải cận sai số nghiệm ODE.

\end{enumerate}

\textbf{Output:} Nghiệm tại các mốc,\allowbreak{} bảng hệ số/\allowbreak{}tầng RK hoặc dự báo–hiệu chỉnh,\allowbreak{} phần dư và thông tin kiểm tra bước ẩn.

\examplebreak\subsubsection*{53. Adams 5}

\textbf{Input:} Vế phải x\ensuremath{\prime};\allowbreak{} y\ensuremath{\prime};\allowbreak{} z\ensuremath{\prime} (\allowbreak{}1–3 hàm);\allowbreak{} Giá trị đầu x,\allowbreak{} y,\allowbreak{} z;\allowbreak{} Thời điểm đầu t0;\allowbreak{} Thời điểm cuối;\allowbreak{} Bước h;\allowbreak{} Dung sai bước ẩn.

\begin{enumerate}

\item Viết hệ u\ensuremath{\prime}=\allowbreak{}f(\allowbreak{}t,\allowbreak{}u),\allowbreak{} u(\allowbreak{}a)=\allowbreak{}u\_0;\allowbreak{} trong ô nhập,\allowbreak{} biến độc lập là t và các trạng thái là x,\allowbreak{}y,\allowbreak{}z.

\item Lập lưới t\_i=\allowbreak{}a+\allowbreak{}ih;\allowbreak{} h phải đúng chiều và kết thúc đúng ở b. Công bố bộ hệ số của phương pháp.

\item Khởi tạo s\ensuremath{-}1 bước bằng RK4;\allowbreak{} từ đủ lịch sử,\allowbreak{} dự báo AB\_s rồi giải phương trình AM\_(\allowbreak{}s+\allowbreak{}1) với dự báo làm điểm đầu.

\item Ghi các tầng/\allowbreak{}dự báo/\allowbreak{}hiệu chỉnh thực sự được tính. Bước ẩn kiểm tra phần dư và Jacobian;\allowbreak{} đánh giá của bộ giải không đồng nghĩa một vòng Newton.

\item Lập bảng nghiệm qua mọi t\_i,\allowbreak{} kèm phần dư bước ẩn. Kết luận trên lưới h;\allowbreak{} dung sai hệ ẩn không phải cận sai số nghiệm ODE.

\end{enumerate}

\textbf{Output:} Nghiệm tại các mốc,\allowbreak{} bảng hệ số/\allowbreak{}tầng RK hoặc dự báo–hiệu chỉnh,\allowbreak{} phần dư và thông tin kiểm tra bước ẩn.

\examplebreak\subsubsection*{54. RK2 tham số}

\textbf{Input:} Vế phải x\ensuremath{\prime};\allowbreak{} y\ensuremath{\prime};\allowbreak{} z\ensuremath{\prime} (\allowbreak{}1–3 hàm);\allowbreak{} Giá trị đầu x,\allowbreak{} y,\allowbreak{} z;\allowbreak{} Thời điểm đầu t0;\allowbreak{} Thời điểm cuối;\allowbreak{} Bước h;\allowbreak{} Dung sai bước ẩn;\allowbreak{} Tham số biết trước;\allowbreak{} Giá trị tham số.

\begin{enumerate}

\item Viết hệ u\ensuremath{\prime}=\allowbreak{}f(\allowbreak{}t,\allowbreak{}u),\allowbreak{} u(\allowbreak{}a)=\allowbreak{}u\_0;\allowbreak{} trong ô nhập,\allowbreak{} biến độc lập là t và các trạng thái là x,\allowbreak{}y,\allowbreak{}z.

\item Lập lưới t\_i=\allowbreak{}a+\allowbreak{}ih;\allowbreak{} h phải đúng chiều và kết thúc đúng ở b. Công bố bộ hệ số của phương pháp.

\item Công bố bảng A,\allowbreak{}b,\allowbreak{}c. Mỗi bước tính lần lượt K\_j=\allowbreak{}f(\allowbreak{}t\_i+\allowbreak{}c\_jh,\allowbreak{}u\_i+\allowbreak{}h\ensuremath{\sum}\_(\allowbreak{}l<j)a\_jl K\_l),\allowbreak{} rồi u\_(\allowbreak{}i+\allowbreak{}1)=\allowbreak{}u\_i+\allowbreak{}h\ensuremath{\sum}\_j b\_jK\_j.

\item Ghi các tầng/\allowbreak{}dự báo/\allowbreak{}hiệu chỉnh thực sự được tính. Bước ẩn kiểm tra phần dư và Jacobian;\allowbreak{} đánh giá của bộ giải không đồng nghĩa một vòng Newton.

\item Lập bảng nghiệm qua mọi t\_i,\allowbreak{} kèm phần dư bước ẩn. Kết luận trên lưới h;\allowbreak{} dung sai hệ ẩn không phải cận sai số nghiệm ODE.

\end{enumerate}

\textbf{Output:} Nghiệm tại các mốc,\allowbreak{} bảng hệ số/\allowbreak{}tầng RK hoặc dự báo–hiệu chỉnh,\allowbreak{} phần dư và thông tin kiểm tra bước ẩn.

\examplebreak\subsubsection*{55. RK3 tham số}

\textbf{Input:} Vế phải x\ensuremath{\prime};\allowbreak{} y\ensuremath{\prime};\allowbreak{} z\ensuremath{\prime} (\allowbreak{}1–3 hàm);\allowbreak{} Giá trị đầu x,\allowbreak{} y,\allowbreak{} z;\allowbreak{} Thời điểm đầu t0;\allowbreak{} Thời điểm cuối;\allowbreak{} Bước h;\allowbreak{} Dung sai bước ẩn;\allowbreak{} alpha2;\allowbreak{} alpha3.

\begin{enumerate}

\item Viết hệ u\ensuremath{\prime}=\allowbreak{}f(\allowbreak{}t,\allowbreak{}u),\allowbreak{} u(\allowbreak{}a)=\allowbreak{}u\_0;\allowbreak{} trong ô nhập,\allowbreak{} biến độc lập là t và các trạng thái là x,\allowbreak{}y,\allowbreak{}z.

\item Lập lưới t\_i=\allowbreak{}a+\allowbreak{}ih;\allowbreak{} h phải đúng chiều và kết thúc đúng ở b. Công bố bộ hệ số của phương pháp.

\item Công bố bảng A,\allowbreak{}b,\allowbreak{}c. Mỗi bước tính lần lượt K\_j=\allowbreak{}f(\allowbreak{}t\_i+\allowbreak{}c\_jh,\allowbreak{}u\_i+\allowbreak{}h\ensuremath{\sum}\_(\allowbreak{}l<j)a\_jl K\_l),\allowbreak{} rồi u\_(\allowbreak{}i+\allowbreak{}1)=\allowbreak{}u\_i+\allowbreak{}h\ensuremath{\sum}\_j b\_jK\_j.

\item Ghi các tầng/\allowbreak{}dự báo/\allowbreak{}hiệu chỉnh thực sự được tính. Bước ẩn kiểm tra phần dư và Jacobian;\allowbreak{} đánh giá của bộ giải không đồng nghĩa một vòng Newton.

\item Lập bảng nghiệm qua mọi t\_i,\allowbreak{} kèm phần dư bước ẩn. Kết luận trên lưới h;\allowbreak{} dung sai hệ ẩn không phải cận sai số nghiệm ODE.

\end{enumerate}

\textbf{Output:} Nghiệm tại các mốc,\allowbreak{} bảng hệ số/\allowbreak{}tầng RK hoặc dự báo–hiệu chỉnh,\allowbreak{} phần dư và thông tin kiểm tra bước ẩn.

\examplebreak\subsubsection*{56. RK tùy chỉnh}

\textbf{Input:} Vế phải x\ensuremath{\prime};\allowbreak{} y\ensuremath{\prime};\allowbreak{} z\ensuremath{\prime} (\allowbreak{}1–3 hàm);\allowbreak{} Giá trị đầu x,\allowbreak{} y,\allowbreak{} z;\allowbreak{} Thời điểm đầu t0;\allowbreak{} Thời điểm cuối;\allowbreak{} Bước h;\allowbreak{} Dung sai bước ẩn;\allowbreak{} Ma trận A (\allowbreak{}JSON);\allowbreak{} Trọng số b;\allowbreak{} Vị trí c.

\begin{enumerate}

\item Viết hệ u\ensuremath{\prime}=\allowbreak{}f(\allowbreak{}t,\allowbreak{}u),\allowbreak{} u(\allowbreak{}a)=\allowbreak{}u\_0;\allowbreak{} trong ô nhập,\allowbreak{} biến độc lập là t và các trạng thái là x,\allowbreak{}y,\allowbreak{}z.

\item Lập lưới t\_i=\allowbreak{}a+\allowbreak{}ih;\allowbreak{} h phải đúng chiều và kết thúc đúng ở b. Công bố bộ hệ số của phương pháp.

\item Công bố bảng A,\allowbreak{}b,\allowbreak{}c. Mỗi bước tính lần lượt K\_j=\allowbreak{}f(\allowbreak{}t\_i+\allowbreak{}c\_jh,\allowbreak{}u\_i+\allowbreak{}h\ensuremath{\sum}\_(\allowbreak{}l<j)a\_jl K\_l),\allowbreak{} rồi u\_(\allowbreak{}i+\allowbreak{}1)=\allowbreak{}u\_i+\allowbreak{}h\ensuremath{\sum}\_j b\_jK\_j.

\item Ghi các tầng/\allowbreak{}dự báo/\allowbreak{}hiệu chỉnh thực sự được tính. Bước ẩn kiểm tra phần dư và Jacobian;\allowbreak{} đánh giá của bộ giải không đồng nghĩa một vòng Newton.

\item Lập bảng nghiệm qua mọi t\_i,\allowbreak{} kèm phần dư bước ẩn. Kết luận trên lưới h;\allowbreak{} dung sai hệ ẩn không phải cận sai số nghiệm ODE.

\end{enumerate}

\textbf{Output:} Nghiệm tại các mốc,\allowbreak{} bảng hệ số/\allowbreak{}tầng RK hoặc dự báo–hiệu chỉnh,\allowbreak{} phần dư và thông tin kiểm tra bước ẩn.

\subsection{Bài toán biên}

\examplebreak\subsubsection*{57. Biên tổng quát}

\textbf{Input:} Hàm p(\allowbreak{}x);\allowbreak{} Hàm q(\allowbreak{}x);\allowbreak{} Hàm f(\allowbreak{}x);\allowbreak{} Đầu đoạn a;\allowbreak{} Cuối đoạn b;\allowbreak{} Bước h;\allowbreak{} Biên trái: alpha beta gamma;\allowbreak{} Biên phải: alpha beta gamma.

\begin{enumerate}

\item Viết \ensuremath{-}(\allowbreak{}p u\ensuremath{\prime})\ensuremath{\prime}+\allowbreak{}qu=\allowbreak{}f trên [a,\allowbreak{}b];\allowbreak{} công bố lưới và quy ước dấu.

\item Tính p tại nửa bước,\allowbreak{} q tại mốc;\allowbreak{} lập các hàng sai phân ở mốc trong từ thông lượng hai mặt ô.

\item Thay \ensuremath{\alpha}u\ensuremath{\prime}+\allowbreak{}\ensuremath{\beta}u=\allowbreak{}\ensuremath{\gamma} ở mỗi đầu;\allowbreak{} trường hợp \ensuremath{\alpha}=\allowbreak{}0 là Dirichlet,\allowbreak{} trường hợp \ensuremath{\alpha}\ensuremath{\ne}0 dùng phương trình nửa ô.

\item Công bố các đường chéo và vế phải;\allowbreak{} giải hệ băng đã lập,\allowbreak{} giữ đủ các hàng biên.

\item Lập bảng u\_i,\allowbreak{} thế lại hệ và báo phần dư. Phần dư đại số không đồng nghĩa sai số so với nghiệm liên tục.

\end{enumerate}

\textbf{Output:} Hệ rời rạc,\allowbreak{} nghiệm trên lưới hoặc các cặp trị riêng–vectơ riêng và thông tin chuẩn hóa/\allowbreak{}kiểm tra.

\examplebreak\subsubsection*{58. Trị riêng}

\textbf{Input:} Hàm p(\allowbreak{}x);\allowbreak{} Hàm q(\allowbreak{}x);\allowbreak{} Hàm r(\allowbreak{}x);\allowbreak{} Đầu đoạn a;\allowbreak{} Cuối đoạn b;\allowbreak{} Bước h;\allowbreak{} Số trị riêng nhỏ nhất.

\begin{enumerate}

\item Viết \ensuremath{-}(\allowbreak{}p u\ensuremath{\prime})\ensuremath{\prime}+\allowbreak{}qu=\allowbreak{}\ensuremath{\lambda}r u trên [a,\allowbreak{}b];\allowbreak{} công bố lưới và quy ước dấu.

\item Tính p tại nửa bước,\allowbreak{} q tại mốc;\allowbreak{} lập các hàng sai phân ở mốc trong từ thông lượng hai mặt ô.

\item Dùng u(\allowbreak{}a)=\allowbreak{}u(\allowbreak{}b)=\allowbreak{}0;\allowbreak{} kiểm tra p>0,\allowbreak{} r>0 và dựng K v=\allowbreak{}\ensuremath{\lambda} R v trên các mốc trong.

\item Đổi sang ma trận đối xứng R\textasciicircum{}(\allowbreak{}\ensuremath{-}1/\allowbreak{}2) K R\textasciicircum{}(\allowbreak{}\ensuremath{-}1/\allowbreak{}2),\allowbreak{} tính các trị riêng nhỏ nhất và đổi lại vectơ.

\item Khôi phục hai biên 0,\allowbreak{} chuẩn hóa max|u|=\allowbreak{}1 đúng như tool;\allowbreak{} công bố từng trị riêng và vectơ riêng trên lưới.

\end{enumerate}

\textbf{Output:} Hệ rời rạc,\allowbreak{} nghiệm trên lưới hoặc các cặp trị riêng–vectơ riêng và thông tin chuẩn hóa/\allowbreak{}kiểm tra.

